\chapter{Technical derivations for Graphical Polyhedral Algebra}
\label{appendix:gpa-proofs}

Chapter~\ref{sec:gpa} states the normal forms and comparison facts of the
existing theory, and retains the soundness, expressivity, and completeness
proofs in the main text. The auxiliary order, elimination, and normal-form
derivations are collected here. The full theory, including comparison of
generating matrices and the empty case, is established in
\cite{bonchi2021polyhedral}.

\section{The order structure of the inequality generator}\label{sec:gpa-rules}

The central computational primitive for reasoning inside this algebra is
Fourier--Motzkin Elimination, which projects a polyhedral cone onto a
lower-dimensional space by eliminating one variable at a time.

Rescaling and sign are given outright by (P5) and (P6), and the pairwise merge
by (spider). What the axioms do not state, and what we need first, is that
$\inlinegen{ge}$ deserves its name: reflexivity and transitivity are not among
them, and antisymmetry is postulated only as the inclusion
$\subseteq$.

\begin{lemma}[The generator is an order]\label{lem:order}
    In $\mathbf{GPA}_{\rtop}$ the generator $\inlinegen{ge}$ is reflexive,
    transitive and antisymmetric.
\end{lemma}

\noindent The proof is collected in Appendix~\ref{appendix:gpa-proof-1}.

Since $\inlinegen{coge}$ is $\inlinegen{ge}$ with an antipode on each side, by
(P6) at $k=-1$, every axiom comes in a mirrored version, and we use both
without further comment.

The move the axioms leave implicit is the transport of a term across a sum: it
is the step that turns a row of a system into a bound on a single variable,
and it is the content of the next lemma.

\begin{lemma}[Cap form]\label{lem:capform}
    In $\mathbf{GPA}_{\rtop}$,
    % An inequality may be moved across a white sum; and a row s+u >= 0 is the
% cap form of s >= -u.
\begin{center}
    \begin{tikzpicture}[axpic]
        \node at (-1.0,0.5) {$(C1)$};
        \pic (g1) at (0.500,1.000) {ge};
        \pic (g2) at (1.750,0.500) {multiplication};
        \coordinate (id_in_1)  at (0.000,0.000);
        \coordinate (id_out_1) at (1.250,0.000);
        \draw (id_in_1) -- (id_out_1);
        \draw (g1-out-0) to[out=0,in=180] (g2-in-0);
        \draw (id_out_1) to[out=0,in=180] (g2-in-1);
    \end{tikzpicture}
    \;\begin{tikzpicture}[axpic]\node{$=$};\end{tikzpicture}\;
    \begin{tikzpicture}[axpic]
        \pic (g1) at (0.500,0.500) {multiplication};
        \pic (g2) at (1.750,0.500) {ge};
        \draw (g1-out-0) to[out=0,in=180] (g2-in-0);
    \end{tikzpicture}
    \;\begin{tikzpicture}[axpic]\node{$=$};\end{tikzpicture}\;
    \begin{tikzpicture}[axpic]
        \pic (g1) at (0.500,0.000) {ge};
        \pic (g2) at (1.750,0.500) {multiplication};
        \coordinate (id_in_1)  at (0.000,1.000);
        \coordinate (id_out_1) at (1.250,1.000);
        \draw (id_in_1) -- (id_out_1);
        \draw (id_out_1) to[out=0,in=180] (g2-in-0);
        \draw (g1-out-0) to[out=0,in=180] (g2-in-1);
    \end{tikzpicture} \\
    \vspace{4mm}

    \begin{tikzpicture}[axpic]
        \node at (-1.0,0.5) {$(C2)$};
        \pic (g1) at (0.500,0.500) {multiplication};
        \pic (g2) at (1.750,0.500) {ge};
        \pic (g3) at (3.000,0.500) {cozero};
        \draw (g1-out-0) to[out=0,in=180] (g2-in-0);
        \draw (g2-out-0) to[out=0,in=180] (g3-in-0);
    \end{tikzpicture}
    \;\begin{tikzpicture}[axpic]\node{$=$};\end{tikzpicture}\;
    \begin{tikzpicture}[axpic]
        \pic (g1) at (0.500,1.000) {ge};
        \pic (g2) at (0.500,0.000) {antipode};
        \draw (g1-out-0) -- (1.600,1.000) arc (90:-90:0.500) -- (g2-out-0);
    \end{tikzpicture}
\end{center}

\end{lemma}

The first chain reads $x+u\ge v$ iff $v-u\le x$: an inequality may be applied
before or after a sum with the same result. The second reads $s+u\ge0$ iff
$s\ge-u$, and it is the form we shall use: the diagram
$(\inlinegen{ge}\oplus\mathrm{id});\varepsilon$, whose denotation is
$\{(\ell,r)\mid \ell\ge r\}$, we call the \emph{cap form} of $\ell\ge r$, so
that a row $s+u\ge0$ of a system is the cap form of the bound $s\ge-u$.

\noindent The proof is collected in Appendix~\ref{appendix:gpa-proof-2}.

The same bending turns a two-sided system $Ax\ge By$ into a homogeneous one.
This is the presentation in which elimination takes place, since it has a
single block of rows and no distinguished side.

\begin{lemma}[Cone form]\label{lem:coneform}
    Let $A:n\to p$ and $B:m\to p$ be matrices. Bending the $m$ outputs of
    $A;\inlinegen{ge};B^{\mathrm{op}}$ into inputs yields the cone
    $[A\;\;{-B}];\inlinegen{ge}^{\oplus p};\mathsf{coz}^{\oplus p}$, and
    unbending recovers the original diagram.
    % Cone form: bending the outputs of a two-sided normal form turns Ax >= By
% into the homogeneous cone Ax - By >= 0.
\begin{center}
    \begin{tikzpicture}[axpic]
        \pic[val=A] (g1) at (0.500,0.500) {matrix};
        \pic (g2) at (1.750,0.500) {ge};
        \pic[val=B] (g3) at (3.000,0.500) {comatrix};
        \draw (g1-out-0) to[out=0,in=180] (g2-in-0);
        \draw (g2-out-0) to[out=0,in=180] (g3-in-0);
        \draw (g3-out-0) -- (4.100,0.500) arc (90:-90:0.500) -- (0.000,-0.500);
        \node[font=\scriptsize] at (0.050,0.700) {$n$};
        \node[font=\scriptsize] at (2.375,0.700) {$p$};
        \node[font=\scriptsize] at (0.150,-0.300) {$m$};
    \end{tikzpicture}
    \;\begin{tikzpicture}[axpic]\node{$=$};\end{tikzpicture}\;
    \begin{tikzpicture}[axpic]
        \pic[val={[A \ {-B}]}] (g1) at (1.100,0.000) {2dmatrix};
        \pic (g2) at (2.600,0.000) {ge};
        \pic (g3) at (3.850,0.000) {cozero};
        \coordinate (id_in_1) at (0.000,0.300);
        \coordinate (id_in_2) at (0.000,-0.300);
        \draw (id_in_1) -- (g1-in-0);
        \draw (id_in_2) -- (g1-in-1);
        \draw (g1-out-0) to[out=0,in=180] (g2-in-0);
        \draw (g2-out-0) to[out=0,in=180] (g3-in-0);
        \node[font=\scriptsize] at (0.050,0.500) {$n$};
        \node[font=\scriptsize] at (0.050,-0.500) {$m$};
        \node[font=\scriptsize] at (2.150,0.200) {$p$};
    \end{tikzpicture}
\end{center}

\end{lemma}

That is, $\exists v.\;Ax\ge v,\ v=By$ iff $Ax-By\ge0$. Rescaling the rows of a
cone by positive constants leaves it unchanged, since $Dz\ge0$ iff $z\ge0$ for
diagonal $D>0$; this is (P5) again, applied row by row.

\noindent The proof is collected in Appendix~\ref{appendix:gpa-proof-3}.

The third move, the pairwise merge, is (spider). It is the only one of the
three that is not local: on its left-hand side a black node forces a single
intermediate value below every $x_i$ and above every $y_j$, and on its
right-hand side that value is gone and only the relations $x_i\ge y_j$ survive.
The axiom is therefore the elimination of an intermediate variable in the form
discussed in the introduction, and the number of rows grows from $a+b$ to
$ab$ exactly as it does in the classical algorithm.

\section{Fourier--Motzkin Elimination}\label{sec:gpa-fme}

These three transformations allow us to eliminate a single variable. Let
$M\in\mathbb{R}^{p\times q}$ have first column $a$, and split the row indices
into $P=\{i\mid a_i>0\}$, $N=\{j\mid a_j<0\}$ and $Z=\{i\mid a_i=0\}$.

\begin{proposition}[Fourier-Motzkin Elimination]\label{prop:fme}
    The following diagrammatic procedure implements the FME algorithm:
    \begin{figure}[H]
        \centering
        \diageq{
  \pic (g1) at (0.500,0.500) {ge};
  \begin{scope}[shift={(1.250,0.000)}]
  \pic[val=A] (g2) at (0.500,0.500) {comatrix};
  \end{scope}
  \draw (g1-out-0) to[out=0,in=180] (g2-in-0);
  \begin{scope}[shift={(2.500,0.000)}]
  \pic[val=B] (g3) at (0.500,0.500) {matrix};
  \end{scope}
  \draw (g2-out-0) to[out=0,in=180] (g3-in-0);
  \begin{scope}[shift={(3.750,0.000)}]
  \pic (g4) at (0.500,0.500) {ge};
  \end{scope}
  \draw (g3-out-0) to[out=0,in=180] (g4-in-0);
      }{
    \pic [val=A'] (g1) at (0.500,0.500) {matrix};
    \begin{scope}[shift={(1.250,0.000)}]
    \pic (g2) at (0.500,0.500) {ge};
    \end{scope}
    \draw (g1-out-0) to[out=0,in=180] (g2-in-0);
    \begin{scope}[shift={(2.500,0.000)}]
    \pic [val=B'](g3) at (0.500,0.500) {comatrix};
    \end{scope}
    \draw (g2-out-0) to[out=0,in=180] (g3-in-0);
}

    \end{figure}
\end{proposition}

\noindent The proof is collected in Appendix~\ref{appendix:gpa-proof-4}.
A pair $(x,y)$ is related by $\inlinegen{ge};g;\inlinegen{ge}$ precisely when
there are $v,w$ with $x\ge v$, $(v,w)\in\llbracket g\rrbracket$ and $w\ge y$,
so the two bound variables $v,w$ are what the elimination has to remove.


\section{Polarity as a technical construction}

The finitely generated normal form is stated directly in
Proposition~\ref{prop:gpa-generated-form}. The source derives it using
polarity; this construction is not needed to read the main presentation.

\begin{definition}[Polar Operator]
    The \emph{polar operator} $-^\circ$ is the involution on generators
    defined by:
    \begin{figure}[H]
        \centering
        \begin{center}\begin{tabular}{c c c @{\hspace{1cm}} c c c}
  % copy |-> comult
  \begin{tikzpicture}[axpic]
      \pic (a) at (0.500,0.500) {copy};
      \draw (-0.250,0.500) -- (a-in-0);
      \draw (a-out-0) -- (1.200,1.000);
      \draw (a-out-1) -- (1.200,0.000);
  \end{tikzpicture} &
  \begin{tikzpicture}[axpic]\node{$\mapsto$};\end{tikzpicture} &
  \begin{tikzpicture}[axpic]
      \pic (a) at (0.500,0.500) {comultiplication};
      \draw (-0.250,0.500) -- (a-in-0);
      \draw (a-out-0) -- (1.200,1.000);
      \draw (a-out-1) -- (1.200,0.000);
  \end{tikzpicture} &

  % cocopy |-> mult
  \begin{tikzpicture}[axpic]
      \pic (a) at (0.500,0.500) {cocopy};
      \draw (-0.250,1.000) -- (a-in-0);
      \draw (-0.250,0.000) -- (a-in-1);
      \draw (a-out-0) -- (1.200,0.500);
  \end{tikzpicture} &
  \begin{tikzpicture}[axpic]\node{$\mapsto$};\end{tikzpicture} &
  \begin{tikzpicture}[axpic]
      \pic (a) at (0.500,0.500) {multiplication};
      \draw (-0.250,1.000) -- (a-in-0);
      \draw (-0.250,0.000) -- (a-in-1);
      \draw (a-out-0) -- (1.200,0.500);
  \end{tikzpicture} \\[4mm]

  % discard |-> cozero
  \begin{tikzpicture}[axpic]
      \pic (a) at (0.500,0.500) {discard};
      \draw (-0.250,0.500) -- (a-in-0);
  \end{tikzpicture} &
  \begin{tikzpicture}[axpic]\node{$\mapsto$};\end{tikzpicture} &
  \begin{tikzpicture}[axpic]
      \pic (a) at (0.500,0.500) {cozero};
      \draw (-0.250,0.500) -- (a-in-0);
  \end{tikzpicture} &

  % codiscard |-> zero
  \begin{tikzpicture}[axpic]
      \pic (a) at (0.500,0.500) {codiscard};
      \draw (a-out-0) -- (1.200,0.500);
  \end{tikzpicture} &
  \begin{tikzpicture}[axpic]\node{$\mapsto$};\end{tikzpicture} &
  \begin{tikzpicture}[axpic]
      \pic (a) at (0.500,0.500) {zero};
      \draw (a-out-0) -- (1.200,0.500);
  \end{tikzpicture} \\[4mm]

  % scalar(k) |-> coscalar(k)
  \begin{tikzpicture}[axpic]
      \pic[val=k] (a) at (0.500,0.500) {scalar};
      \draw (-0.250,0.500) -- (a-in-0);
      \draw (a-out-0) -- (1.200,0.500);
  \end{tikzpicture} &
  \begin{tikzpicture}[axpic]\node{$\mapsto$};\end{tikzpicture} &
  \begin{tikzpicture}[axpic]
      \pic[val=k] (a) at (0.500,0.500) {coscalar};
      \draw (-0.250,0.500) -- (a-in-0);
      \draw (a-out-0) -- (1.200,0.500);
  \end{tikzpicture} &

  % coscalar(k) |-> scalar(k)
  \begin{tikzpicture}[axpic]
      \pic[val=k] (a) at (0.500,0.500) {coscalar};
      \draw (-0.250,0.500) -- (a-in-0);
      \draw (a-out-0) -- (1.200,0.500);
  \end{tikzpicture} &
  \begin{tikzpicture}[axpic]\node{$\mapsto$};\end{tikzpicture} &
  \begin{tikzpicture}[axpic]
      \pic[val=k] (a) at (0.500,0.500) {scalar};
      \draw (-0.250,0.500) -- (a-in-0);
      \draw (a-out-0) -- (1.200,0.500);
  \end{tikzpicture} \\[4mm]

  % mult |-> cocopy
  \begin{tikzpicture}[axpic]
      \pic (a) at (0.500,0.500) {multiplication};
      \draw (-0.250,1.000) -- (a-in-0);
      \draw (-0.250,0.000) -- (a-in-1);
      \draw (a-out-0) -- (1.200,0.500);
  \end{tikzpicture} &
  \begin{tikzpicture}[axpic]\node{$\mapsto$};\end{tikzpicture} &
  \begin{tikzpicture}[axpic]
      \pic (a) at (0.500,0.500) {cocopy};
      \draw (-0.250,1.000) -- (a-in-0);
      \draw (-0.250,0.000) -- (a-in-1);
      \draw (a-out-0) -- (1.200,0.500);
  \end{tikzpicture} &

  % comult |-> copy
  \begin{tikzpicture}[axpic]
      \pic (a) at (0.500,0.500) {comultiplication};
      \draw (-0.250,0.500) -- (a-in-0);
      \draw (a-out-0) -- (1.200,1.000);
      \draw (a-out-1) -- (1.200,0.000);
  \end{tikzpicture} &
  \begin{tikzpicture}[axpic]\node{$\mapsto$};\end{tikzpicture} &
  \begin{tikzpicture}[axpic]
      \pic (a) at (0.500,0.500) {copy};
      \draw (-0.250,0.500) -- (a-in-0);
      \draw (a-out-0) -- (1.200,1.000);
      \draw (a-out-1) -- (1.200,0.000);
  \end{tikzpicture} \\[4mm]

  % zero |-> codiscard
  \begin{tikzpicture}[axpic]
      \pic (a) at (0.500,0.500) {zero};
      \draw (a-out-0) -- (1.200,0.500);
  \end{tikzpicture} &
  \begin{tikzpicture}[axpic]\node{$\mapsto$};\end{tikzpicture} &
  \begin{tikzpicture}[axpic]
      \pic (a) at (0.500,0.500) {codiscard};
      \draw (a-out-0) -- (1.200,0.500);
  \end{tikzpicture} &

  % cozero |-> discard
  \begin{tikzpicture}[axpic]
      \pic (a) at (0.500,0.500) {cozero};
      \draw (-0.250,0.500) -- (a-in-0);
  \end{tikzpicture} &
  \begin{tikzpicture}[axpic]\node{$\mapsto$};\end{tikzpicture} &
  \begin{tikzpicture}[axpic]
      \pic (a) at (0.500,0.500) {discard};
      \draw (-0.250,0.500) -- (a-in-0);
  \end{tikzpicture} \\
\end{tabular}\end{center}

% ==================== second row: ge |-> copy;((ge;cozero)(x)id) ====================
\begin{center}\begin{tabular}{c c c}
  \begin{tikzpicture}[axpic]
      \pic (g) at (0.500,0.500) {ge};
      \draw (-0.250,0.500) -- (g-in-0);
      \draw (g-out-0) -- (1.200,0.500);
  \end{tikzpicture}
  &
  \begin{tikzpicture}[axpic]\node{$\mapsto$};\end{tikzpicture}
  &
  \begin{tikzpicture}[axpic]
      \pic (cp) at (0.500,0.500) {copy};
      \pic (g)  at (1.750,1.000) {ge};
      \pic (cz) at (3.000,1.000) {cozero};
      \draw (-0.250,0.500) -- (cp-in-0);
      \draw (cp-out-0) to[out=0,in=180] (g-in-0);
      \draw (g-out-0) to[out=0,in=180] (cz-in-0);
      \draw (cp-out-1) -- (3.750,0.000);
  \end{tikzpicture}
\end{tabular}\end{center}

    \end{figure}
    Its action on diagrams is defined functorially by
    $(a \circ b)^\circ = a^\circ \circ b^\circ$.
\end{definition}

\section{The order generator}\label{appendix:gpa-proof-1}

\begin{proof}
    Transitivity is (spider) with one input fed by the black unit and one
    output discarded, the plugs being absorbed by (P3) and the black
    (co)unit laws:
    % Transitivity of the inequality generator: the (spider) axiom with one input
% fed by the black unit and one output discarded.
% ge;ge = (ge * (codiscard;ge); cocopy ; copy; (ge ; (ge; discard) = (id(1) * codiscard); (cocopy * cocopy); (ge * ge); (id(1) * swap(1,1) *id(1)); (copy * copy); (id(1) * discard)
\begin{center}
    \begin{tikzpicture}[axpic]
        \pic (g1) at (0.500,0.500) {ge};
        \pic (g2) at (1.750,0.500) {ge};
        \draw (g1-out-0) to[out=0,in=180] (g2-in-0);
    \end{tikzpicture} \\
    \vspace{3mm}

    \;\begin{tikzpicture}[axpic]\node{$\overset{(\text{unit})}{=}$};\end{tikzpicture}\;
    \begin{tikzpicture}[axpic]
        \pic (u)  at (0.400,0.000) {codiscard};
        \pic (g1) at (1.900,1.000) {ge};
        \pic (g2) at (1.900,0.000) {ge};
        \pic (c1) at (3.150,0.500) {cocopy};
        \pic (c2) at (4.400,0.500) {copy};
        \pic (h1) at (5.650,1.000) {ge};
        \pic (h2) at (5.650,0.000) {ge};
        \pic (d)  at (6.900,0.000) {discard};
        \draw (0.900,1.000) -- (g1-in-0);
        \draw (u-out-0) -- (g2-in-0);
        \draw (g1-out-0) to[out=0,in=180] (c1-in-0);
        \draw (g2-out-0) to[out=0,in=180] (c1-in-1);
        \draw (c1-out-0) -- (c2-in-0);
        \draw (c2-out-0) to[out=0,in=180] (h1-in-0);
        \draw (c2-out-1) to[out=0,in=180] (h2-in-0);
        \draw (h1-out-0) -- (6.900,1.000);
        \draw (h2-out-0) -- (d-in-0);
    \end{tikzpicture} \\
    \vspace{3mm}

    \;\begin{tikzpicture}[axpic]\node{$\overset{(\text{spider})}{=}$};\end{tikzpicture}\;
    \begin{tikzpicture}[axpic]
        \pic (u)   at (0.400,-0.500) {codiscard};
        \pic (c1)  at (1.400,1.500) {copy};
        \pic (c2)  at (1.400,-0.500) {copy};
        \pic (g11) at (2.900,2.000) {ge};
        \pic (g12) at (2.900,1.000) {ge};
        \pic (g21) at (2.900,0.000) {ge};
        \pic (g22) at (2.900,-1.000) {ge};
        \pic (d1)  at (4.650,1.500) {cocopy};
        \pic (d2)  at (4.650,-0.500) {cocopy};
        \pic (dd)  at (5.900,-0.500) {discard};
        \draw (0.400,1.500) -- (c1-in-0);
        \draw (u-out-0) -- (c2-in-0);
        \draw (c1-out-0) to[out=0,in=180] (g11-in-0);
        \draw (c1-out-1) to[out=0,in=180] (g12-in-0);
        \draw (c2-out-0) to[out=0,in=180] (g21-in-0);
        \draw (c2-out-1) to[out=0,in=180] (g22-in-0);
        \draw (g11-out-0) to[out=0,in=180] (d1-in-0);
        \draw (g21-out-0) to[out=0,in=180] (d1-in-1);
        \draw (g12-out-0) to[out=0,in=180] (d2-in-0);
        \draw (g22-out-0) to[out=0,in=180] (d2-in-1);
        \draw (d1-out-0) -- (5.900,1.500);
        \draw (d2-out-0) -- (dd-in-0);
    \end{tikzpicture} \\
    \vspace{3mm}

    \;\begin{tikzpicture}[axpic]\node{$\overset{(P3),(\text{unit})}{=}$};\end{tikzpicture}\;
    \begin{tikzpicture}[axpic]
        \pic (g1) at (0.500,0.500) {ge};
    \end{tikzpicture}
\end{center}


    \noindent For reflexivity, we first establish an auxiliary identity, denoted by $H1$:
    \begin{center}
        % ===================== TOP ROW: derivation of (*) =====================
\begin{center}\begin{tabular}{c c c}
  % Panel A: comult;(ge (x) id)
  \begin{tikzpicture}[axpic]
      \node at (-0.5, 0.5) {$(H1)$};
      \pic (c) at (0.500,0.500) {comultiplication};
      \pic (g) at (1.750,1.000) {ge};
      \draw (c-out-0) to[out=0,in=180] (g-in-0);
      \draw (c-out-1) to[out=0,in=180] (2.500,0.000);
  \end{tikzpicture}
  \begin{tikzpicture}[axpic]\node{$\overset{ IH}{=}$};\end{tikzpicture} &

  % Panel B: (id (x) (codiscard;copy)); (id (x) antipode (x) id); ((mult;ge) (x) id)
  \begin{tikzpicture}[axpic]
      \pic (cd) at (0.500,-0.500) {codiscard};
      \pic (cp) at (1.750,-0.500) {copy};
      \pic (ap) at (3.500, 0.000) {antipode};
      \pic (m)  at (5.000, 0.500) {multiplication};
      \pic (g)  at (6.750, 0.500) {ge};
      \draw (-0.250,1.000) -- (m-in-0);
      \draw (cd-out-0) to[out=0,in=180] (cp-in-0);
      \draw (cp-out-0) to[out=0,in=180] (ap-in-0);
      \draw (ap-out-0) to[out=0,in=180] (m-in-1);
      \draw (cp-out-1) -- (7.500,-1.000);
      \draw (m-out-0) to[out=0,in=180] (g-in-0);
      \draw (g-out-0) -- (7.500, 0.500);
  \end{tikzpicture}\\
  \begin{tikzpicture}[axpic]\node{$\overset{ P2}{=}$};\end{tikzpicture} &

  % Panel C: (ge_A) (x) (antipode;ge)_B  ; mult (x) id
  \begin{tikzpicture}[axpic]
      \pic (ga) at (0.500, 1.000) {ge};
      \pic (cd) at (0.500,-0.500) {codiscard};
      \pic (cp) at (1.750,-0.500) {copy};
      \pic (ap) at (3.000, 0.000) {antipode};
      \pic (gb) at (4.250, 0.000) {ge};
      \pic (m)  at (5.750, 0.500) {multiplication};
      \draw (-0.250,1.000) -- (ga-in-0);
      \draw (ga-out-0) to[out=0,in=180] (m-in-0);
      \draw (cd-out-0) to[out=0,in=180] (cp-in-0);
      \draw (cp-out-0) to[out=0,in=180] (ap-in-0);
      \draw (ap-out-0) -- (gb-in-0);
      \draw (gb-out-0) to[out=0,in=180] (m-in-1);
      \draw (cp-out-1) -- (6.500,-1.000);
  \end{tikzpicture}
  \\ \begin{tikzpicture}[axpic]\node{$\overset{ P6}{=}$};\end{tikzpicture}&

  % Panel D: (ge_A) (x) (coge;antipode)_B ; mult (x) id
  \begin{tikzpicture}[axpic]
      \pic (ga) at (0.500, 1.000) {ge};
      \pic (cd) at (0.500,-0.500) {codiscard};
      \pic (cp) at (1.750,-0.500) {copy};
      \pic (cg) at (3.000, 0.000) {coge};
      \pic (ap) at (4.250, 0.000) {antipode};
      \pic (m)  at (5.750, 0.500) {multiplication};
      \draw (-0.250,1.000) -- (ga-in-0);
      \draw (ga-out-0) to[out=0,in=180] (m-in-0);
      \draw (cd-out-0) to[out=0,in=180] (cp-in-0);
      \draw (cp-out-0) to[out=0,in=180] (cg-in-0);
      \draw (cg-out-0) -- (ap-in-0);
      \draw (ap-out-0) to[out=0,in=180] (m-in-1);
      \draw (cp-out-1) -- (6.500,-1.000);
  \end{tikzpicture}
  \\[2mm]
\end{tabular}\end{center}

\begin{center}\begin{tabular}{c c c c}
  \begin{tikzpicture}[axpic]\node{$\overset{ IH}{=}$};\end{tikzpicture} &

  % Panel E = (*):  ge;comult;(id (x) ge)
  \begin{tikzpicture}[axpic]
      \pic (g1) at (0.500,0.500) {ge};
      \pic (c)  at (1.750,0.500) {comultiplication};
      \pic (g2) at (3.000,0.000) {ge};
      \draw (-0.250,0.500) -- (g1-in-0);
      \draw (g1-out-0) to[out=0,in=180] (c-in-0);
      \draw (c-out-0) -- (3.750,1.000);
      \draw (c-out-1) to[out=0,in=180] (g2-in-0);
      \draw (g2-out-0) -- (3.750,0.000);
  \end{tikzpicture}
\end{tabular}\end{center}

% ===================== BOTTOM ROW: reflexivity of >= =====================
\begin{center}\begin{tabular}{c c c c c c c c c c c}
  % id
  \begin{tikzpicture}[axpic]
      \draw (0.000,0.500) -- (1.500,0.500);
  \end{tikzpicture}
  & \begin{tikzpicture}[axpic]\node{$\overset{ IH}{=}$};\end{tikzpicture} &

  % comult (open dangling outputs)
  \begin{tikzpicture}[axpic]
      \pic (c) at (0.500,0.500) {comultiplication};
      \draw (-0.250,0.500) -- (c-in-0);
      \draw (c-out-0) -- (1.500,1.000);
      \draw (c-out-1) -- (1.500,0.000);
      \pic at (1.500,1.000) {cozero};
  \end{tikzpicture}
  & \begin{tikzpicture}[axpic]\node{$\overset{ P4}{\subseteq}$};\end{tikzpicture} &

  % comult;(ge (x) id)  (open dangling outputs)
  \begin{tikzpicture}[axpic]
      \pic (c) at (0.500,0.500) {comultiplication};
      \pic (g) at (1.750,1.000) {ge};
      \draw (-0.250,0.500) -- (c-in-0);
      \draw (c-out-0) to[out=0,in=180] (g-in-0);
      \draw (g-out-0) -- (2.500,1.000);
      \draw (c-out-1) -- (2.500,0.000);
      \pic at (2.500,1.000) {cozero};
  \end{tikzpicture}
  & \begin{tikzpicture}[axpic]\node{$\overset{ (H1)}{=}$};\end{tikzpicture} &

  % ge;comult;(id (x) ge)  (open dangling outputs)
  \begin{tikzpicture}[axpic]
      \pic (g1) at (0.500,0.500) {ge};
      \pic (c)  at (1.750,0.500) {comultiplication};
      \pic (g2) at (3.000,0.000) {ge};
      \draw (-0.250,0.500) -- (g1-in-0);
      \draw (g1-out-0) to[out=0,in=180] (c-in-0);
      \draw (c-out-0) -- (3.750,1.000);
      \draw (c-out-1) to[out=0,in=180] (g2-in-0);
      \draw (g2-out-0) -- (3.750,0.000);
      \pic at (3.750,1.000) {cozero};
  \end{tikzpicture}
  \\[2mm]
\end{tabular}\end{center}

\begin{center}\begin{tabular}{c c c c c}
  \begin{tikzpicture}[axpic]\node{$\overset{ IH}{=}$};\end{tikzpicture} &

  % ge;ge
  \begin{tikzpicture}[axpic]
      \pic (g1) at (0.500,0.500) {ge};
      \pic (g2) at (1.750,0.500) {ge};
      \draw (-0.250,0.500) -- (g1-in-0);
      \draw (g1-out-0) to[out=0,in=180] (g2-in-0);
  \end{tikzpicture}
  & \begin{tikzpicture}[axpic]\node{$\overset{ trans}{=}$};\end{tikzpicture} &

  % ge
  \begin{tikzpicture}[axpic]
      \pic (g) at (0.500,0.500) {ge};
      \draw (-0.250,0.500) -- (g-in-0);
  \end{tikzpicture}
\end{tabular}\end{center}

    \end{center}

    \noindent Antisymmetry, the converse of the inclusion (antisym):
    \begin{center}
        \begin{tikzpicture}
    \node{(1)};
\end{tikzpicture}\;
\begin{tikzpicture}[axpic]
  \begin{scope}[shift={(0.000,0.500)}]
  \pic (g1) at (0.500,0.500) {copy};
  \end{scope}
  \begin{scope}[shift={(1.250,0.000)}]
  \begin{scope}[shift={(0.000,1.000)}]
  \pic (g2) at (0.500,0.500) {ge};
  \end{scope}
  \pic (g3) at (0.500,0.500) {coge};
  \end{scope}
  \draw (g1-out-0) to[out=0,in=180] (g2-in-0);
  \draw (g1-out-1) to[out=0,in=180] (g3-in-0);
  \begin{scope}[shift={(2.500,0.500)}]
  \pic (g4) at (0.500,0.500) {cocopy};
  \end{scope}
  \draw (g2-out-0) to[out=0,in=180] (g4-in-0);
  \draw (g3-out-0) to[out=0,in=180] (g4-in-1);
\end{tikzpicture}
\;\begin{tikzpicture}[axpic]\node{$\subseteq$};\end{tikzpicture}\;
\begin{tikzpicture}[axpic]
  \coordinate (id_in_1) at (0,0.500);
  \coordinate (id_out_2) at (1.000,0.500);
  \draw (id_in_1) -- (id_out_2);
\end{tikzpicture}

\vspace{5mm}

\begin{tikzpicture}
    \node{(2)};
\end{tikzpicture}\;
\begin{tikzpicture}[axpic]
  \coordinate (id_in_1) at (0,0.500);
  \coordinate (id_out_2) at (1.000,0.500);
  \draw (id_in_1) -- (id_out_2);
\end{tikzpicture}
\;\begin{tikzpicture}[axpic]\node{$=$};\end{tikzpicture}\;
\begin{tikzpicture}[axpic]
  \pic (g1) at (0.500,0.500) {copy};
  \begin{scope}[shift={(1.250,0.000)}]
  \pic (g2) at (0.500,0.500) {cocopy};
  \end{scope}
  \draw (g1-out-0) to[out=0,in=180] (g2-in-0);
  \draw (g1-out-1) to[out=0,in=180] (g2-in-1);
\end{tikzpicture}
\;\begin{tikzpicture}[axpic]\node{$\subseteq$};\end{tikzpicture}\;
\begin{tikzpicture}[axpic]
  \begin{scope}[shift={(0.000,0.500)}]
  \pic (g1) at (0.500,0.500) {copy};
  \end{scope}
  \begin{scope}[shift={(1.250,0.000)}]
  \begin{scope}[shift={(0.000,1.000)}]
  \pic (g2) at (0.500,0.500) {ge};
  \end{scope}
  \pic (g3) at (0.500,0.500) {coge};
  \end{scope}
  \draw (g1-out-0) to[out=0,in=180] (g2-in-0);
  \draw (g1-out-1) to[out=0,in=180] (g3-in-0);
  \begin{scope}[shift={(2.500,0.500)}]
  \pic (g4) at (0.500,0.500) {cocopy};
  \end{scope}
  \draw (g2-out-0) to[out=0,in=180] (g4-in-0);
  \draw (g3-out-0) to[out=0,in=180] (g4-in-1);
\end{tikzpicture}

\vspace{5mm}
\;\begin{tikzpicture}[axpic]\node{$(1)\wedge(2)\implies$};\end{tikzpicture}\;
\begin{tikzpicture}[axpic]
  \coordinate (id_in_1) at (0,0.500);
  \coordinate (id_out_2) at (1.000,0.500);
  \draw (id_in_1) -- (id_out_2);
\end{tikzpicture}
\;\begin{tikzpicture}[axpic]\node{$=$};\end{tikzpicture}\;
\begin{tikzpicture}[axpic]
  \begin{scope}[shift={(0.000,0.500)}]
  \pic (g1) at (0.500,0.500) {copy};
  \end{scope}
  \begin{scope}[shift={(1.250,0.000)}]
  \begin{scope}[shift={(0.000,1.000)}]
  \pic (g2) at (0.500,0.500) {ge};
  \end{scope}
  \pic (g3) at (0.500,0.500) {coge};
  \end{scope}
  \draw (g1-out-0) to[out=0,in=180] (g2-in-0);
  \draw (g1-out-1) to[out=0,in=180] (g3-in-0);
  \begin{scope}[shift={(2.500,0.500)}]
  \pic (g4) at (0.500,0.500) {cocopy};
  \end{scope}
  \draw (g2-out-0) to[out=0,in=180] (g4-in-0);
  \draw (g3-out-0) to[out=0,in=180] (g4-in-1);
\end{tikzpicture}

    \end{center}
\end{proof}

\section{Cap form}\label{appendix:gpa-proof-2}

\begin{proof}
    Bending the second input of the white sum to an output is a bijection on
    hom-sets, so the first equation may be compared in bent form; the third
    follows by commutativity of the sum. Closing the bent leg with a white
    cozero then gives the second equation.
    % Proof of (C1): an inequality may be applied before or after a white sum.
% Proof of (C2): closing the sum with a white cozero gives the cap form.
% Broken over lines to fit the text block; all diagrams share the
% convention that the sum sits at y=1.0, its top leg at y=1.5 and its bottom
% leg at y=0.5, so that legs meet the sum without kinks.
\begin{center}
    \begin{tikzpicture}[axpic]\node{$(C1)$};\end{tikzpicture}\;
    \begin{tikzpicture}[axpic]
        \pic (g1) at (0.500,1.500) {ge};
        \pic (a1) at (1.750,1.000) {multiplication};
        \coordinate (id_in_1)  at (0.000,0.500);
        \coordinate (id_out_1) at (1.000,0.500);
        \draw (id_in_1) -- (id_out_1);
        \draw (g1-out-0) -- (a1-in-0);
        \draw (id_out_1) -- (a1-in-1);
        \draw (a1-out-0) -- (2.750,1.000);
    \end{tikzpicture}
    \;\begin{tikzpicture}[axpic]\node{$=$};\end{tikzpicture}\;
    \begin{tikzpicture}[axpic]
        \pic (g1) at (0.500,1.500) {ge};
        \pic (g2) at (1.750,1.500) {ge};
        \pic (a1) at (3.000,1.000) {multiplication};
        \coordinate (id_in_1)  at (0.000,0.500);
        \coordinate (id_out_1) at (2.250,0.500);
        \draw (id_in_1) -- (id_out_1);
        \draw (g1-out-0) -- (g2-in-0);
        \draw (g2-out-0) -- (a1-in-0);
        \draw (id_out_1) -- (a1-in-1);
        \draw (a1-out-0) -- (4.000,1.000);
    \end{tikzpicture}
    \;\begin{tikzpicture}[axpic]\node{$=$};\end{tikzpicture}\;
    \begin{tikzpicture}[axpic]
        \pic (g1) at (0.500,1.500) {ge};
        \pic (g2) at (0.500,0.500) {ge};
        \pic (a1) at (1.750,1.000) {multiplication};
        \pic (g3) at (3.000,1.000) {ge};
        \draw (g1-out-0) -- (a1-in-0);
        \draw (g2-out-0) -- (a1-in-1);
        \draw (a1-out-0) -- (g3-in-0);
        \draw (g3-out-0) -- (3.500,1.000);
    \end{tikzpicture} \\
    \vspace{3mm}

    \;\begin{tikzpicture}[axpic]\node{$=$};\end{tikzpicture}\;
    \begin{tikzpicture}[axpic]
        \pic (a1) at (0.750,1.000) {multiplication};
        \pic (g1) at (2.000,1.000) {ge};
        \pic (g2) at (3.250,1.000) {ge};
        \draw (a1-out-0) -- (g1-in-0);
        \draw (g1-out-0) -- (g2-in-0);
        \draw (g2-out-0) -- (3.750,1.000);
    \end{tikzpicture}
    \;\begin{tikzpicture}[axpic]\node{$=$};\end{tikzpicture}\;
    \begin{tikzpicture}[axpic]
        \pic (a1) at (0.750,1.000) {multiplication};
        \pic (g1) at (2.000,1.000) {ge};
        \draw (a1-out-0) -- (g1-in-0);
        \draw (g1-out-0) -- (2.500,1.000);
    \end{tikzpicture}
\end{center}

\vspace{4mm}

\begin{center}
    \begin{tikzpicture}[axpic]\node{$(C2)$};\end{tikzpicture}\;
    \begin{tikzpicture}[axpic]
        \pic (a1) at (0.750,1.000) {multiplication};
        \pic (g1) at (2.000,1.000) {ge};
        \pic (z1) at (3.250,1.000) {cozero};
        \draw (a1-out-0) -- (g1-in-0);
        \draw (g1-out-0) -- (z1-in-0);
    \end{tikzpicture}
    \;\begin{tikzpicture}[axpic]\node{$\overset{(C1)}{=}$};\end{tikzpicture}\;
    \begin{tikzpicture}[axpic]
        \pic (g1) at (0.500,1.500) {ge};
        \pic (a1) at (1.750,1.000) {multiplication};
        \pic (z1) at (3.000,1.000) {cozero};
        \coordinate (id_in_1)  at (0.000,0.500);
        \coordinate (id_out_1) at (1.000,0.500);
        \draw (id_in_1) -- (id_out_1);
        \draw (g1-out-0) -- (a1-in-0);
        \draw (id_out_1) -- (a1-in-1);
        \draw (a1-out-0) -- (z1-in-0);
    \end{tikzpicture}
    \;\begin{tikzpicture}[axpic]\node{$\overset{(\text{white cap})}{=}$};\end{tikzpicture}\;
    \begin{tikzpicture}[axpic]
        \pic (g1) at (0.500,1.500) {ge};
        \pic (s1) at (1.750,1.500) {antipode};
        \pic (c1) at (3.000,1.000) {cocopy};
        \pic (d1) at (4.250,1.000) {discard};
        \coordinate (id_in_1)  at (0.000,0.500);
        \coordinate (id_out_1) at (2.250,0.500);
        \draw (id_in_1) -- (id_out_1);
        \draw (g1-out-0) -- (s1-in-0);
        \draw (s1-out-0) -- (c1-in-0);
        \draw (id_out_1) -- (c1-in-1);
        \draw (c1-out-0) -- (d1-in-0);
    \end{tikzpicture}
\end{center}

\end{proof}

\section{Cone form}\label{appendix:gpa-proof-3}

\begin{proof}
    Bending the matrix $B$ exposes a white sum with an antipode on one leg, and
    the cap form pushes the inequality past it.
    % Proof of the cone form: bending the comatrix produces a white sum with an
% antipode, and the cap-form lemma pushes the inequality past it.
\begin{center}
    \begin{tikzpicture}[axpic]
        \pic[val=A] (g1) at (0.500,0.500) {matrix};
        \pic (g2) at (1.750,0.500) {ge};
        \pic[val=B] (g3) at (3.000,0.500) {comatrix};
        \draw (g1-out-0) to[out=0,in=180] (g2-in-0);
        \draw (g2-out-0) to[out=0,in=180] (g3-in-0);
        \draw (g3-out-0) -- (4.100,0.500) arc (90:-90:0.500) -- (0.000,-0.500);
    \end{tikzpicture}
    \;\begin{tikzpicture}[axpic]\node{$\overset{\text{(white cap)}}{=}$};\end{tikzpicture}\;
    \begin{tikzpicture}[axpic]
        \pic[val=A] (g1) at (0.500,1.000) {matrix};
        \pic (g2) at (1.750,1.000) {ge};
        \pic[val=B] (g3) at (0.500,0.000) {matrix};
        \pic (g4) at (1.750,0.000) {antipode};
        \pic (g5) at (3.000,0.500) {multiplication};
        \pic (g6) at (4.250,0.500) {cozero};
        \draw (g1-out-0) to[out=0,in=180] (g2-in-0);
        \draw (g3-out-0) to[out=0,in=180] (g4-in-0);
        \draw (g2-out-0) to[out=0,in=180] (g5-in-0);
        \draw (g4-out-0) to[out=0,in=180] (g5-in-1);
        \draw (g5-out-0) to[out=0,in=180] (g6-in-0);
    \end{tikzpicture} \\
    \vspace{3mm}

    \;\begin{tikzpicture}[axpic]\node{$\overset{(\text{cap form})}{=}$};\end{tikzpicture}\;
    \begin{tikzpicture}[axpic]
        \pic[val=A] (g1) at (0.500,1.000) {matrix};
        \pic[val=B] (g2) at (0.500,0.000) {matrix};
        \pic (g3) at (1.750,0.000) {antipode};
        \pic (g4) at (3.000,0.500) {multiplication};
        \pic (g5) at (4.250,0.500) {ge};
        \pic (g6) at (5.500,0.500) {cozero};
        \draw (g2-out-0) to[out=0,in=180] (g3-in-0);
        \draw (g1-out-0) to[out=0,in=180] (g4-in-0);
        \draw (g3-out-0) to[out=0,in=180] (g4-in-1);
        \draw (g4-out-0) to[out=0,in=180] (g5-in-0);
        \draw (g5-out-0) to[out=0,in=180] (g6-in-0);
    \end{tikzpicture} \\
    \vspace{3mm}

    \;\begin{tikzpicture}[axpic]\node{$\overset{(\text{G1})}{=}$};\end{tikzpicture}\;
    \begin{tikzpicture}[axpic]
        \pic[val={[A \ {-B}]}] (g1) at (1.100,0.000) {2dmatrix};
        \pic (g2) at (2.600,0.000) {ge};
        \pic (g3) at (3.850,0.000) {cozero};
        \coordinate (id_in_1) at (0.000,0.300);
        \coordinate (id_in_2) at (0.000,-0.300);
        \draw (id_in_1) -- (g1-in-0);
        \draw (id_in_2) -- (g1-in-1);
        \draw (g1-out-0) to[out=0,in=180] (g2-in-0);
        \draw (g2-out-0) to[out=0,in=180] (g3-in-0);
    \end{tikzpicture}
\end{center}

\end{proof}

\section{Fourier--Motzkin elimination}\label{appendix:gpa-proof-4}

\begin{proof}
        \tikzset{bridge/.style={preaction={draw=white,line width=2.6pt}},
 every picture/.append style={line width=.5pt,line cap=round,line join=round}}
First consider one middle wire; ellipses denote parallel repetitions.


\[
\fits{\begin{tikzpicture}[x=.77cm,y=.64cm,scale=1,baseline=(current bounding box.center)]
        \pic[] (g) at (0,0) {ge};
        \pic[val={A},] (a) at (1.95,0) {comatrix};
        \pic[val={B},] (b) at (4.55,0) {matrix};
        \pic[] (h) at (6.5,0) {ge};
        \draw (g-out-0) -- (a-in-0);
        \draw (a-out-0) -- (b-in-0);
        \draw (b-out-0) -- (h-in-0);

    \end{tikzpicture}
\eqby{(matrix def.)}
\begin{tikzpicture}[x=.77cm,y=.64cm,scale=1,baseline=(current bounding box.center)]
        \pic[] (gi0) at (0,1) {ge};
        \pic[val={a_1},] (a0) at (1.35,1) {coscalar};
        \pic[val={b_1},] (b0) at (5.15,1) {scalar};
        \pic[] (go0) at (6.5,1) {ge};
        \draw (gi0-out-0) -- (a0-in-0);
        \draw (b0-out-0) -- (go0-in-0);
        \pic[] (gi1) at (0,-1) {ge};
        \pic[val={a_n},] (a1) at (1.35,-1) {coscalar};
        \pic[val={b_m},] (b1) at (5.15,-1) {scalar};
        \pic[] (go1) at (6.5,-1) {ge};
        \draw (gi1-out-0) -- (a1-in-0);
        \draw (b1-out-0) -- (go1-in-0);
        \pic[yscale=2] (mu) at (2.65,0) {cocopy};
        \pic[yscale=2] (de) at (3.85,0) {copy};
        \draw (mu-out-0) -- (de-in-0);
        \draw (a0-out-0) -- (mu-in-0);
        \draw (de-out-0) -- (b0-in-0);
        \draw (a1-out-0) -- (mu-in-1);
        \draw (de-out-1) -- (b1-in-0);
        \node[font=\scriptsize,anchor=center] at (0,0) {$\vdots$};
        \node[font=\scriptsize,anchor=center] at (6.5,0) {$\vdots$};

\end{tikzpicture}}
\]

For positive coefficients,


\[
\fits{\begin{tikzpicture}[x=.77cm,y=.64cm,scale=1,baseline=(current bounding box.center)]
\pic[] (a) at (0,0) {ge};
\pic[val={d},] (b) at (1.3,0) {coscalar};
\begin{pgfonlayer}{background}\draw (a-out-0) -- (b-in-0);\end{pgfonlayer}
\end{tikzpicture}\eqby{inversion lemma}\begin{tikzpicture}[x=.77cm,y=.64cm,scale=1,baseline=(current bounding box.center)]
\pic[] (a) at (0,0) {ge};
\pic[val={d^{-1}},] (b) at (1.3,0) {scalar};
\begin{pgfonlayer}{background}\draw (a-out-0) -- (b-in-0);\end{pgfonlayer}
\end{tikzpicture}\eqby{(P5)}\begin{tikzpicture}[x=.77cm,y=.64cm,scale=1,baseline=(current bounding box.center)]
\pic[val={d^{-1}},] (a) at (0,0) {scalar};
\pic[] (b) at (1.3,0) {ge};
\begin{pgfonlayer}{background}\draw (a-out-0) -- (b-in-0);\end{pgfonlayer}
\end{tikzpicture}\eqby{inversion lemma}\begin{tikzpicture}[x=.77cm,y=.64cm,scale=1,baseline=(current bounding box.center)]
\pic[val={d},] (a) at (0,0) {coscalar};
\pic[] (b) at (1.3,0) {ge};
\begin{pgfonlayer}{background}\draw (a-out-0) -- (b-in-0);\end{pgfonlayer}
\end{tikzpicture}}
\]

\[
\fits{\begin{tikzpicture}[x=.77cm,y=.64cm,scale=1,baseline=(current bounding box.center)]
            \pic[val={d},] (n0) at (0.5,0) {scalar};
            \pic[] (n1) at (1.85,0) {ge};
            \draw (n0-out-0) -- (n1-in-0);

            \end{tikzpicture}
\eqby{(P5)}
\begin{tikzpicture}[x=.77cm,y=.64cm,scale=1,baseline=(current bounding box.center)]
            \pic[] (n0) at (0.5,0) {ge};
            \pic[val={d},] (n1) at (1.85,0) {scalar};
            \draw (n0-out-0) -- (n1-in-0);

        \end{tikzpicture}}
\]

and for negative coefficients,


\[
\fits{\begin{tikzpicture}[x=.77cm,y=.64cm,scale=1,baseline=(current bounding box.center)]
\pic[] (a) at (0,0) {ge};
\pic[val={d},] (b) at (1.3,0) {coscalar};
\begin{pgfonlayer}{background}\draw (a-out-0) -- (b-in-0);\end{pgfonlayer}
\end{tikzpicture}\eqby{inversion lemma}\begin{tikzpicture}[x=.77cm,y=.64cm,scale=1,baseline=(current bounding box.center)]
\pic[] (a) at (0,0) {ge};
\pic[val={d^{-1}},] (b) at (1.3,0) {scalar};
\begin{pgfonlayer}{background}\draw (a-out-0) -- (b-in-0);\end{pgfonlayer}
\end{tikzpicture}\eqby{(P6)}\begin{tikzpicture}[x=.77cm,y=.64cm,scale=1,baseline=(current bounding box.center)]
\pic[val={d^{-1}},] (a) at (0,0) {scalar};
\pic[] (b) at (1.3,0) {coge};
\begin{pgfonlayer}{background}\draw (a-out-0) -- (b-in-0);\end{pgfonlayer}
\end{tikzpicture}\eqby{inversion lemma}\begin{tikzpicture}[x=.77cm,y=.64cm,scale=1,baseline=(current bounding box.center)]
\pic[val={d},] (a) at (0,0) {coscalar};
\pic[] (b) at (1.3,0) {coge};
\begin{pgfonlayer}{background}\draw (a-out-0) -- (b-in-0);\end{pgfonlayer}
\end{tikzpicture}}
\]

\[
\fits{\begin{tikzpicture}[x=.77cm,y=.64cm,scale=1,baseline=(current bounding box.center)]
            \pic[val={d},] (n0) at (0.5,0) {scalar};
            \pic[] (n1) at (1.85,0) {ge};
            \draw (n0-out-0) -- (n1-in-0);

            \end{tikzpicture}
\eqby{(P6)}
\begin{tikzpicture}[x=.77cm,y=.64cm,scale=1,baseline=(current bounding box.center)]
            \pic[] (n0) at (0.5,0) {coge};
            \pic[val={d},] (n1) at (1.85,0) {scalar};
            \draw (n0-out-0) -- (n1-in-0);

        \end{tikzpicture}}
\]

Zero legs detach:


\[
\fits{\begin{tikzpicture}[x=.77cm,y=.64cm,scale=1,baseline=(current bounding box.center)]
        \pic[] (n0) at (0.5,0) {ge};
        \pic[val={0},] (n1) at (1.85,0) {coscalar};
        \draw (n0-out-0) -- (n1-in-0);

        \end{tikzpicture}
\eqby{(zero sc.), (slide)}
\begin{tikzpicture}[x=.77cm,y=.64cm,scale=1,baseline=(current bounding box.center)]
        \pic[] (g) at (0.5,0) {ge};
        \pic[] (z) at (1.75,0) {cozero};
        \pic[] (d) at (2.95,0) {codiscard};
        \draw (g-out-0) -- (z-in-0);

    \end{tikzpicture}}
\]

\[
\fits{\begin{tikzpicture}[x=.77cm,y=.64cm,scale=1,baseline=(current bounding box.center)]
        \pic[val={0},] (n0) at (0.5,0) {scalar};
        \pic[] (n1) at (1.85,0) {ge};
        \draw (n0-out-0) -- (n1-in-0);

        \end{tikzpicture}
\eqby{(zero sc.)}
\begin{tikzpicture}[x=.77cm,y=.64cm,scale=1,baseline=(current bounding box.center)]
        \pic[] (d) at (0.5,0) {discard};
        \pic[] (z) at (1.65,0) {zero};
        \pic[] (g) at (2.85,0) {ge};
        \draw (z-out-0) -- (g-in-0);

    \end{tikzpicture}}
\]


\[
\fits{\begin{tikzpicture}[x=.77cm,y=.64cm,scale=1,baseline=(current bounding box.center)]
\pic[] (u) at (0,-0.5) {codiscard};
\pic[] (m) at (1.4,0) {cocopy};
\begin{pgfonlayer}{background}\draw (-.5,.5) -- (m-in-0);\end{pgfonlayer}
\begin{pgfonlayer}{background}\draw (u-out-0) -- (m-in-1);\end{pgfonlayer}
\end{tikzpicture}
\eqby{(unit)}
\begin{tikzpicture}[x=.77cm,y=.64cm,scale=1,baseline=(current bounding box.center)]
\begin{pgfonlayer}{background}\draw (0,0) -- (1.4,0);\end{pgfonlayer}
\end{tikzpicture}\qquad\begin{tikzpicture}[x=.77cm,y=.64cm,scale=1,baseline=(current bounding box.center)]
\pic[] (m) at (0,0) {copy};
\pic[] (u) at (1.4,-0.5) {discard};
\begin{pgfonlayer}{background}\draw (m-out-0) -- (1.9,.5);\end{pgfonlayer}
\begin{pgfonlayer}{background}\draw (m-out-1) -- (u-in-0);\end{pgfonlayer}
\end{tikzpicture}
\eqby{(unit)}
\begin{tikzpicture}[x=.77cm,y=.64cm,scale=1,baseline=(current bounding box.center)]
\begin{pgfonlayer}{background}\draw (0,0) -- (1.4,0);\end{pgfonlayer}
\end{tikzpicture}}
\]

Bend the reversed legs:


\[
\fits{\begin{tikzpicture}[x=.77cm,y=.64cm,scale=1,baseline=(current bounding box.center)]
\pic[] (g) at (0.5,0.65) {coge};
\begin{pgfonlayer}{background}\draw (g-out-0) -- (1.55,.65) to[out=0,in=0,looseness=1.4] (1.55,-.65) -- (0,-.65);\end{pgfonlayer}
\end{tikzpicture}
\eqby{(slide)}
\begin{tikzpicture}[x=.77cm,y=.64cm,scale=1,baseline=(current bounding box.center)]
\pic[] (g) at (0.5,-0.65) {ge};
\begin{pgfonlayer}{background}\draw (0,.65) -- (1.55,.65) to[out=0,in=0,looseness=1.4] (1.55,-.65) -- (g-out-0);\end{pgfonlayer}
\end{tikzpicture}}
\]



The nonzero part now rewrites as follows.


\[
\fits{\begin{tikzpicture}[x=.77cm,y=.64cm,scale=1,baseline=(current bounding box.center)]
        \pic[val={d_1},] (a0) at (0,1) {coscalar};
        \pic[val={e_1},] (b0) at (7.2,1) {scalar};
        \pic[val={d_p},] (a1) at (0,-1) {coscalar};
        \pic[val={e_q},] (b1) at (7.2,-1) {scalar};
        \pic[] (g0) at (1.5,1) {ge};
        \pic[] (h0) at (5.7,1) {ge};
        \draw (a0-out-0) -- (g0-in-0);
        \draw (h0-out-0) -- (b0-in-0);
        \pic[] (g1) at (1.5,-1) {ge};
        \pic[] (h1) at (5.7,-1) {ge};
        \draw (a1-out-0) -- (g1-in-0);
        \draw (h1-out-0) -- (b1-in-0);
        \pic[yscale=2] (mu) at (2.85,0) {cocopy};
        \pic[yscale=2] (de) at (4.35,0) {copy};
        \draw (mu-out-0) -- (de-in-0);
        \draw (g0-out-0) -- (mu-in-0);
        \draw (de-out-0) -- (h0-in-0);
        \draw (g1-out-0) -- (mu-in-1);
        \draw (de-out-1) -- (h1-in-0);
        \node[font=\scriptsize,anchor=center] at (0,0) {$\vdots$};
        \node[font=\scriptsize,anchor=center] at (7.2,0) {$\vdots$};

    \end{tikzpicture}
\eqby{generalised (spider)}
\begin{tikzpicture}[x=.77cm,y=.64cm,scale=1,baseline=(current bounding box.center)]
        \pic[val={d_1},] (a0) at (0,1) {coscalar};
        \pic[val={e_1},] (b0) at (7.2,1) {scalar};
        \pic[val={d_p},] (a1) at (0,-1) {coscalar};
        \pic[val={e_q},] (b1) at (7.2,-1) {scalar};
        \pic[] (d0) at (1.5,1) {copy};
        \pic[] (m0) at (5.7,1) {cocopy};
        \draw (a0-out-0) -- (d0-in-0);
        \draw (m0-out-0) -- (b0-in-0);
        \pic[] (d1) at (1.5,-1) {copy};
        \pic[] (m1) at (5.7,-1) {cocopy};
        \draw (a1-out-0) -- (d1-in-0);
        \draw (m1-out-0) -- (b1-in-0);
        \pic[] (g0) at (3.1,1.5) {ge};
        \pic[] (g1) at (3.1,0.5) {ge};
        \pic[] (g2) at (3.1,-0.5) {ge};
        \pic[] (g3) at (3.1,-1.5) {ge};
        \draw (d0-out-0) -- (g0-in-0);
        \draw (d0-out-1) -- (g1-in-0);
        \draw (d1-out-0) -- (g2-in-0);
        \draw (d1-out-1) -- (g3-in-0);
        \draw (g0-out-0) -- (m0-in-0);
        \draw (g3-out-0) -- (m1-in-1);
        \draw (g1-out-0) to[out=0,in=180] (m1-in-0);
        \draw[bridge] (g2-out-0) to[out=0,in=180] (m0-in-1);
        \node[font=\scriptsize,anchor=center] at (0,0) {$\vdots$};
        \node[font=\scriptsize,anchor=center] at (7.2,0) {$\vdots$};

\node[font=\scriptsize] at (3.1,1) {$\vdots$};
\node[font=\scriptsize] at (3.1,0) {$\vdots$};
\node[font=\scriptsize] at (3.1,-1) {$\vdots$};
\end{tikzpicture}}
\]


\[
\fits{\eqby{inversion lemma}\quad\begin{tikzpicture}[x=.77cm,y=.64cm,scale=1,baseline=(current bounding box.center)]
        \pic[val={d_1^{-1}},] (a0) at (0,1) {scalar};
        \pic[val={e_1^{-1}},] (b0) at (7.2,1) {coscalar};
        \pic[val={d_p^{-1}},] (a1) at (0,-1) {scalar};
        \pic[val={e_q^{-1}},] (b1) at (7.2,-1) {coscalar};
        \pic[] (d0) at (1.5,1) {copy};
        \pic[] (m0) at (5.7,1) {cocopy};
        \draw (a0-out-0) -- (d0-in-0);
        \draw (m0-out-0) -- (b0-in-0);
        \pic[] (d1) at (1.5,-1) {copy};
        \pic[] (m1) at (5.7,-1) {cocopy};
        \draw (a1-out-0) -- (d1-in-0);
        \draw (m1-out-0) -- (b1-in-0);
        \pic[] (g0) at (3.1,1.5) {ge};
        \pic[] (g1) at (3.1,0.5) {ge};
        \pic[] (g2) at (3.1,-0.5) {ge};
        \pic[] (g3) at (3.1,-1.5) {ge};
        \draw (d0-out-0) -- (g0-in-0);
        \draw (d0-out-1) -- (g1-in-0);
        \draw (d1-out-0) -- (g2-in-0);
        \draw (d1-out-1) -- (g3-in-0);
        \draw (g0-out-0) -- (m0-in-0);
        \draw (g3-out-0) -- (m1-in-1);
        \draw (g1-out-0) to[out=0,in=180] (m1-in-0);
        \draw[bridge] (g2-out-0) to[out=0,in=180] (m0-in-1);
        \node[font=\scriptsize,anchor=center] at (0,0) {$\vdots$};
        \node[font=\scriptsize,anchor=center] at (7.2,0) {$\vdots$};

\node[font=\scriptsize] at (3.1,1) {$\vdots$};
\node[font=\scriptsize] at (3.1,0) {$\vdots$};
\node[font=\scriptsize] at (3.1,-1) {$\vdots$};
\end{tikzpicture}}
\]


\[
\fits{\eqby{(matrix def.)}\quad\begin{tikzpicture}[x=.77cm,y=.64cm,scale=1,baseline=(current bounding box.center)]
        \pic[val={H},] (a) at (0,0) {matrix};
        \pic[] (g) at (2,0) {ge};
        \pic[val={K},] (b) at (4,0) {comatrix};
        \draw (a-out-0) -- (g-in-0);
        \draw (g-out-0) -- (b-in-0);

\end{tikzpicture}}
\]

\needspace{28mm}
For an empty family of bounds,


\[
\fits{\begin{tikzpicture}[x=.77cm,y=.64cm,scale=1,baseline=(current bounding box.center)]
\pic[] (g) at (0,0.55) {ge};
\pic[] (h) at (0,-0.55) {ge};
\pic[yscale=1.1] (m) at (1.4,0) {cocopy};
\pic[] (e) at (2.7,0) {discard};
\begin{pgfonlayer}{background}\draw (g-out-0) -- (m-in-0);\end{pgfonlayer}
\begin{pgfonlayer}{background}\draw (h-out-0) -- (m-in-1);\end{pgfonlayer}
\begin{pgfonlayer}{background}\draw (m-out-0) -- (e-in-0);\end{pgfonlayer}
\node[font=\scriptsize] at (0,0) {$\vdots$};
\end{tikzpicture}
\eqby{generalised (spider)}
\begin{tikzpicture}[x=.77cm,y=.64cm,scale=1,baseline=(current bounding box.center)]
\pic[] (a) at (0,0.55) {discard};
\pic[] (b) at (0,-0.55) {discard};
\node[font=\scriptsize] at (0,0) {$\vdots$};
\end{tikzpicture}}
\]


\[
\fits{\begin{tikzpicture}[x=.77cm,y=.64cm,scale=1,baseline=(current bounding box.center)]
\pic[] (s) at (0,0) {codiscard};
\pic[yscale=1.1] (d) at (1.3,0) {copy};
\pic[] (g) at (2.7,0.55) {ge};
\pic[] (h) at (2.7,-0.55) {ge};
\begin{pgfonlayer}{background}\draw (s-out-0) -- (d-in-0);\end{pgfonlayer}
\begin{pgfonlayer}{background}\draw (d-out-0) -- (g-in-0);\end{pgfonlayer}
\begin{pgfonlayer}{background}\draw (d-out-1) -- (h-in-0);\end{pgfonlayer}
\node[font=\scriptsize] at (2.7,0) {$\vdots$};
\end{tikzpicture}
\eqby{generalised (spider)}
\begin{tikzpicture}[x=.77cm,y=.64cm,scale=1,baseline=(current bounding box.center)]
\pic[] (a) at (0,0.55) {codiscard};
\pic[] (b) at (0,-0.55) {codiscard};
\node[font=\scriptsize] at (0,0) {$\vdots$};
\end{tikzpicture}}
\]


\[
\fits{\begin{tikzpicture}[x=.77cm,y=.64cm,scale=1,baseline=(current bounding box.center)]
\pic[] (s) at (0,0) {codiscard};
\pic[] (e) at (1.2,0) {discard};
\begin{pgfonlayer}{background}\draw (s-out-0) -- (e-in-0);\end{pgfonlayer}
\end{tikzpicture}
\eqby{empty-scalar lemma}
\begin{tikzpicture}[x=.77cm,y=.64cm,scale=1,baseline=(current bounding box.center)]
\path[use as bounding box] (0,-.2) rectangle (.65,.2);\draw[dashed] (.1,-.18) rectangle (.5,.18);
\end{tikzpicture}}
\]

Bend the detached output constraints as well:


\[
\fits{\begin{tikzpicture}[x=.77cm,y=.64cm,scale=1,baseline=(current bounding box.center)]
\pic[] (z) at (0,0.6) {zero};
\pic[] (g) at (1.3,0.6) {ge};
\begin{pgfonlayer}{background}\draw (z-out-0) -- (g-in-0);\end{pgfonlayer}
\begin{pgfonlayer}{background}\draw (g-out-0) -- (2.2,.6) to[out=0,in=0] (2.2,-.6) -- (0,-.6);\end{pgfonlayer}
\end{tikzpicture}
\eqby{cone form}
\begin{tikzpicture}[x=.77cm,y=.64cm,scale=1,baseline=(current bounding box.center)]
\pic[] (s) at (0.3,0) {antipode};
\pic[] (g) at (1.6,0) {ge};
\pic[] (z) at (2.9,0) {cozero};
\begin{pgfonlayer}{background}\draw (s-out-0) -- (g-in-0);\end{pgfonlayer}
\begin{pgfonlayer}{background}\draw (g-out-0) -- (z-in-0);\end{pgfonlayer}
\end{tikzpicture}}
\]



Reattach the zero branches; the permutation restores their boundary positions.


\[
\fits{\begin{tikzpicture}[x=.77cm,y=.64cm,scale=0.88,baseline=(current bounding box.center)]
\begin{scope}[matrixnode_left/.append style={minimum height=26mm}]
\pic[val={\Pi},] (p) at (0,-0.45) {2dmatrix};
\end{scope}
\begin{pgfonlayer}{background}\draw (-2.65,.3) -- (0,.3);\end{pgfonlayer}
\begin{pgfonlayer}{background}\draw (-2.65,-.3) -- (0,-.3);\end{pgfonlayer}
\pic[val={H},] (h) at (3.7,1.4) {matrix};
\pic[] (g) at (5.75,1.4) {ge};
\pic[val={K},] (k) at (7.8,1.4) {comatrix};
\begin{pgfonlayer}{background}\draw (0,1.4) -- (h-in-0);\end{pgfonlayer}
\begin{pgfonlayer}{background}\draw (h-out-0) -- (g-in-0);\end{pgfonlayer}
\begin{pgfonlayer}{background}\draw (g-out-0) -- (k-in-0);\end{pgfonlayer}
\begin{pgfonlayer}{background}\draw (k-out-0) -- (9.1,1.4) to[out=0,in=0,looseness=1.35] (9.1,0) -- (0,0);\end{pgfonlayer}
\pic[] (ga) at (4.9,-1.2) {ge};
\pic[] (za) at (6.2,-1.2) {cozero};
\begin{pgfonlayer}{background}\draw (0,-1.2) -- (ga-in-0);\end{pgfonlayer}
\begin{pgfonlayer}{background}\draw (ga-out-0) -- (za-in-0);\end{pgfonlayer}
\pic[] (sa) at (3.4,-2.3) {antipode};
\pic[] (gb) at (4.9,-2.3) {ge};
\pic[] (zb) at (6.2,-2.3) {cozero};
\begin{pgfonlayer}{background}\draw (0,-2.3) -- (sa-in-0);\end{pgfonlayer}
\begin{pgfonlayer}{background}\draw (sa-out-0) -- (gb-in-0);\end{pgfonlayer}
\begin{pgfonlayer}{background}\draw (gb-out-0) -- (zb-in-0);\end{pgfonlayer}
\end{tikzpicture}}
\]


\[
\fits{\eqby{cone form}\quad\begin{tikzpicture}[x=.77cm,y=.64cm,scale=0.88,baseline=(current bounding box.center)]
\begin{scope}[matrixnode_left/.append style={minimum height=26mm}]
\pic[val={\Pi},] (p) at (0,-0.45) {2dmatrix};
\end{scope}
\begin{pgfonlayer}{background}\draw (-2.65,.3) -- (0,.3);\end{pgfonlayer}
\begin{pgfonlayer}{background}\draw (-2.65,-.3) -- (0,-.3);\end{pgfonlayer}
\pic[val={H},] (h) at (3.7,1.4) {matrix};
\pic[val={K},] (k) at (3.7,0) {matrix};
\begin{pgfonlayer}{background}\draw (0,1.4) -- (h-in-0);\end{pgfonlayer}
\begin{pgfonlayer}{background}\draw (0,0) -- (k-in-0);\end{pgfonlayer}
\pic[] (s) at (5.4,0) {antipode};
\begin{pgfonlayer}{background}\draw (k-out-0) -- (s-in-0);\end{pgfonlayer}
\pic[yscale=1.4] (sum) at (7.5,0.7) {multiplication};
\begin{pgfonlayer}{background}\draw (h-out-0) -- (sum-in-0);\end{pgfonlayer}
\begin{pgfonlayer}{background}\draw (s-out-0) -- (sum-in-1);\end{pgfonlayer}
\pic[] (g) at (8.9,0.7) {ge};
\pic[] (z) at (10.2,0.7) {cozero};
\begin{pgfonlayer}{background}\draw (sum-out-0) -- (g-in-0);\end{pgfonlayer}
\begin{pgfonlayer}{background}\draw (g-out-0) -- (z-in-0);\end{pgfonlayer}
\pic[] (ga) at (4.9,-1.2) {ge};
\pic[] (za) at (6.2,-1.2) {cozero};
\begin{pgfonlayer}{background}\draw (0,-1.2) -- (ga-in-0);\end{pgfonlayer}
\begin{pgfonlayer}{background}\draw (ga-out-0) -- (za-in-0);\end{pgfonlayer}
\pic[] (sa) at (3.4,-2.3) {antipode};
\pic[] (gb) at (4.9,-2.3) {ge};
\pic[] (zb) at (6.2,-2.3) {cozero};
\begin{pgfonlayer}{background}\draw (0,-2.3) -- (sa-in-0);\end{pgfonlayer}
\begin{pgfonlayer}{background}\draw (sa-out-0) -- (gb-in-0);\end{pgfonlayer}
\begin{pgfonlayer}{background}\draw (gb-out-0) -- (zb-in-0);\end{pgfonlayer}
\end{tikzpicture}}
\]


\[
\fits{\eqby{(comp.)}\quad\begin{tikzpicture}[x=.77cm,y=.64cm,scale=0.88,baseline=(current bounding box.center)]
\begin{scope}[matrixnode_left/.append style={minimum height=26mm}]
\pic[val={\Pi},] (p) at (0,-0.45) {2dmatrix};
\end{scope}
\begin{pgfonlayer}{background}\draw (-2.65,.3) -- (0,.3);\end{pgfonlayer}
\begin{pgfonlayer}{background}\draw (-2.65,-.3) -- (0,-.3);\end{pgfonlayer}
\pic[val={H},] (h) at (3.7,1.4) {matrix};
\pic[val={J},] (k) at (3.7,0) {matrix};
\begin{pgfonlayer}{background}\draw (0,1.4) -- (h-in-0);\end{pgfonlayer}
\begin{pgfonlayer}{background}\draw (0,0) -- (k-in-0);\end{pgfonlayer}
\pic[yscale=1.4] (sum) at (7.5,0.7) {multiplication};
\begin{pgfonlayer}{background}\draw (h-out-0) -- (sum-in-0);\end{pgfonlayer}
\begin{pgfonlayer}{background}\draw (k-out-0) -- (sum-in-1);\end{pgfonlayer}
\pic[] (g) at (8.9,0.7) {ge};
\pic[] (z) at (10.2,0.7) {cozero};
\begin{pgfonlayer}{background}\draw (sum-out-0) -- (g-in-0);\end{pgfonlayer}
\begin{pgfonlayer}{background}\draw (g-out-0) -- (z-in-0);\end{pgfonlayer}
\pic[] (ga) at (4.9,-1.2) {ge};
\pic[] (za) at (6.2,-1.2) {cozero};
\begin{pgfonlayer}{background}\draw (0,-1.2) -- (ga-in-0);\end{pgfonlayer}
\begin{pgfonlayer}{background}\draw (ga-out-0) -- (za-in-0);\end{pgfonlayer}
\pic[] (sa) at (3.4,-2.3) {antipode};
\pic[] (gb) at (4.9,-2.3) {ge};
\pic[] (zb) at (6.2,-2.3) {cozero};
\begin{pgfonlayer}{background}\draw (0,-2.3) -- (sa-in-0);\end{pgfonlayer}
\begin{pgfonlayer}{background}\draw (sa-out-0) -- (gb-in-0);\end{pgfonlayer}
\begin{pgfonlayer}{background}\draw (gb-out-0) -- (zb-in-0);\end{pgfonlayer}
\end{tikzpicture}}
\]


\[
\fits{\eqby{(row)}\quad\begin{tikzpicture}[x=.77cm,y=.64cm,scale=0.88,baseline=(current bounding box.center)]
\begin{scope}[matrixnode_left/.append style={minimum height=26mm}]
\pic[val={\Pi},] (p) at (0,-0.45) {2dmatrix};
\end{scope}
\begin{pgfonlayer}{background}\draw (-2.65,.3) -- (0,.3);\end{pgfonlayer}
\begin{pgfonlayer}{background}\draw (-2.65,-.3) -- (0,-.3);\end{pgfonlayer}
\pic[val={L},] (l) at (4.1,0.7) {2dmatrix};
\begin{pgfonlayer}{background}\draw (0,1.4) -- (2.5,1.4) to[out=0,in=180] (l-in-0);\end{pgfonlayer}
\begin{pgfonlayer}{background}\draw (0,0) -- (2.5,0) to[out=0,in=180] (l-in-1);\end{pgfonlayer}
\pic[] (g) at (6.3,0.7) {ge};
\pic[] (z) at (7.6,0.7) {cozero};
\begin{pgfonlayer}{background}\draw (l-out-0) -- (g-in-0);\end{pgfonlayer}
\begin{pgfonlayer}{background}\draw (g-out-0) -- (z-in-0);\end{pgfonlayer}
\pic[] (ga) at (4.9,-1.2) {ge};
\pic[] (za) at (6.2,-1.2) {cozero};
\begin{pgfonlayer}{background}\draw (0,-1.2) -- (ga-in-0);\end{pgfonlayer}
\begin{pgfonlayer}{background}\draw (ga-out-0) -- (za-in-0);\end{pgfonlayer}
\pic[] (sa) at (3.4,-2.3) {antipode};
\pic[] (gb) at (4.9,-2.3) {ge};
\pic[] (zb) at (6.2,-2.3) {cozero};
\begin{pgfonlayer}{background}\draw (0,-2.3) -- (sa-in-0);\end{pgfonlayer}
\begin{pgfonlayer}{background}\draw (sa-out-0) -- (gb-in-0);\end{pgfonlayer}
\begin{pgfonlayer}{background}\draw (gb-out-0) -- (zb-in-0);\end{pgfonlayer}
\end{tikzpicture}}
\]


\[
\fits{\eqby{(matrix def.)}\quad\begin{tikzpicture}[x=.77cm,y=.64cm,scale=1,baseline=(current bounding box.center)]
\pic[val={\Pi},] (p) at (0,0) {2dmatrix};
\pic[val={T},] (m) at (2.5,0) {matrix};
\begin{pgfonlayer}{background}\draw (p-out-0) -- (m-in-0);\end{pgfonlayer}
\pic[] (g) at (4.5,0) {ge};
\pic[] (z) at (5.8,0) {cozero};
\begin{pgfonlayer}{background}\draw (m-out-0) -- (g-in-0);\end{pgfonlayer}
\begin{pgfonlayer}{background}\draw (g-out-0) -- (z-in-0);\end{pgfonlayer}
\end{tikzpicture}}
\]


\[
\fits{\eqby{(comp.)}\quad\begin{tikzpicture}[x=.77cm,y=.64cm,scale=1,baseline=(current bounding box.center)]
\pic[val={C},] (m) at (0,0) {2dmatrix};
\pic[] (g) at (2,0) {ge};
\pic[] (z) at (3.3,0) {cozero};
\begin{pgfonlayer}{background}\draw (m-out-0) -- (g-in-0);\end{pgfonlayer}
\begin{pgfonlayer}{background}\draw (g-out-0) -- (z-in-0);\end{pgfonlayer}
\end{tikzpicture}}
\]



Restore the original input--output boundary:


\[
\fits{\begin{tikzpicture}[x=.77cm,y=.64cm,scale=1,baseline=(current bounding box.center)]
\pic[val={C},] (m) at (0,0) {2dmatrix};
\pic[] (g) at (2,0) {ge};
\pic[] (z) at (3.3,0) {cozero};
\begin{pgfonlayer}{background}\draw (m-out-0) -- (g-in-0);\end{pgfonlayer}
\begin{pgfonlayer}{background}\draw (g-out-0) -- (z-in-0);\end{pgfonlayer}
\end{tikzpicture}
\eqby{(row)}
\begin{tikzpicture}[x=.77cm,y=.64cm,scale=1,baseline=(current bounding box.center)]
\pic[val={A'},] (a) at (0,0.75) {matrix};
\pic[val={Q},] (b) at (0,-0.75) {matrix};
\pic[yscale=1.5] (sum) at (3,0) {multiplication};
\begin{pgfonlayer}{background}\draw (a-out-0) -- (sum-in-0);\end{pgfonlayer}
\begin{pgfonlayer}{background}\draw (b-out-0) -- (sum-in-1);\end{pgfonlayer}
\pic[] (g) at (4.3,0) {ge};
\pic[] (z) at (5.55,0) {cozero};
\begin{pgfonlayer}{background}\draw (sum-out-0) -- (g-in-0);\end{pgfonlayer}
\begin{pgfonlayer}{background}\draw (g-out-0) -- (z-in-0);\end{pgfonlayer}
\end{tikzpicture}}
\]


\[
\fits{\eqby{(antipode), (comp.)}\quad\begin{tikzpicture}[x=.77cm,y=.64cm,scale=1,baseline=(current bounding box.center)]
\pic[val={A'},] (a) at (0,0.75) {matrix};
\pic[val={Q},] (b) at (0,-0.75) {matrix};
\pic[] (s) at (1.7,-0.75) {antipode};
\pic[] (t) at (2.8,-0.75) {antipode};
\pic[yscale=1.5] (sum) at (4.1,0) {multiplication};
\begin{pgfonlayer}{background}\draw (a-out-0) -- (sum-in-0);\end{pgfonlayer}
\begin{pgfonlayer}{background}\draw (b-out-0) -- (s-in-0);\end{pgfonlayer}
\begin{pgfonlayer}{background}\draw (s-out-0) -- (t-in-0);\end{pgfonlayer}
\begin{pgfonlayer}{background}\draw (t-out-0) -- (sum-in-1);\end{pgfonlayer}
\pic[] (g) at (5.4,0) {ge};
\pic[] (z) at (6.65,0) {cozero};
\begin{pgfonlayer}{background}\draw (sum-out-0) -- (g-in-0);\end{pgfonlayer}
\begin{pgfonlayer}{background}\draw (g-out-0) -- (z-in-0);\end{pgfonlayer}
\end{tikzpicture}}
\]


\[
\fits{\eqby{(comp.)}\quad\begin{tikzpicture}[x=.77cm,y=.64cm,scale=1,baseline=(current bounding box.center)]
\pic[val={A'},] (a) at (0,0.75) {matrix};
\pic[val={B'},] (b) at (0,-0.75) {matrix};
\pic[] (s) at (1.9,-0.75) {antipode};
\pic[yscale=1.5] (sum) at (3.3,0) {multiplication};
\begin{pgfonlayer}{background}\draw (a-out-0) -- (sum-in-0);\end{pgfonlayer}
\begin{pgfonlayer}{background}\draw (b-out-0) -- (s-in-0);\end{pgfonlayer}
\begin{pgfonlayer}{background}\draw (s-out-0) -- (sum-in-1);\end{pgfonlayer}
\pic[] (g) at (4.6,0) {ge};
\pic[] (z) at (5.85,0) {cozero};
\begin{pgfonlayer}{background}\draw (sum-out-0) -- (g-in-0);\end{pgfonlayer}
\begin{pgfonlayer}{background}\draw (g-out-0) -- (z-in-0);\end{pgfonlayer}
\end{tikzpicture}}
\]


\[
\fits{\eqby{(C2)}\quad\begin{tikzpicture}[x=.77cm,y=.64cm,scale=1,baseline=(current bounding box.center)]
\pic[val={A'},] (a) at (0,0.75) {matrix};
\pic[val={B'},] (b) at (0,-0.75) {matrix};
\pic[] (g) at (2,0.75) {ge};
\begin{pgfonlayer}{background}\draw (a-out-0) -- (g-in-0);\end{pgfonlayer}
\begin{pgfonlayer}{background}\draw (g-out-0) -- (3.1,.75) to[out=0,in=0,looseness=1.3] (3.1,-.75) -- (b-out-0);\end{pgfonlayer}
\end{tikzpicture}}
\]


\[
\fits{\begin{tikzpicture}[x=.77cm,y=.64cm,scale=1,baseline=(current bounding box.center)]
        \pic[] (g) at (0,0) {ge};
        \pic[val={A},] (a) at (1.95,0) {comatrix};
        \pic[val={B},] (b) at (4.55,0) {matrix};
        \pic[] (h) at (6.5,0) {ge};
        \draw (g-out-0) -- (a-in-0);
        \draw (a-out-0) -- (b-in-0);
        \draw (b-out-0) -- (h-in-0);

    \end{tikzpicture}
\eqby{unbend}
\begin{tikzpicture}[x=.77cm,y=.64cm,scale=1,baseline=(current bounding box.center)]
        \pic[val={A'},] (a) at (0,0) {matrix};
        \pic[] (g) at (2,0) {ge};
        \pic[val={B'},] (b) at (4,0) {comatrix};
        \draw (a-out-0) -- (g-in-0);
        \draw (g-out-0) -- (b-in-0);

    \end{tikzpicture}}
\]

Now bend an arbitrary middle bundle and expose its independent sources.


\[
\fits{\begin{tikzpicture}[x=.77cm,y=.64cm,scale=1,baseline=(current bounding box.center)]
        \pic[] (g) at (0,0) {ge};
        \pic[val={A},] (a) at (1.95,0) {comatrix};
        \pic[val={B},] (b) at (4.55,0) {matrix};
        \pic[] (h) at (6.5,0) {ge};
        \draw (g-out-0) -- (a-in-0);
        \draw (a-out-0) -- (b-in-0);
        \draw (b-out-0) -- (h-in-0);
        \draw(h-out-0) to[out=0,in=0,looseness=1.4] (7,-1.15) -- (-.5,-1.15);

    \end{tikzpicture}}
\]


\[
\fits{\eqby{(slide)}\quad\begin{tikzpicture}[x=.77cm,y=.64cm,scale=0.88,baseline=(current bounding box.center)]
\pic[] (u) at (0,0.85) {codiscard};
\pic[yscale=.8] (du) at (1.3,0.85) {copy};
\begin{pgfonlayer}{background}\draw (u-out-0) -- (du-in-0);\end{pgfonlayer}
\pic[] (v) at (0,-0.85) {codiscard};
\pic[yscale=.8] (dv) at (1.3,-0.85) {copy};
\begin{pgfonlayer}{background}\draw (v-out-0) -- (dv-in-0);\end{pgfonlayer}
\node[font=\scriptsize] at (0,0) {$\vdots$};
\begin{scope}[matrixnode_left/.append style={minimum height=10mm}]
\pic[val={A},] (a) at (4,1.35) {2dmatrix};
\end{scope}
\begin{scope}[matrixnode_left/.append style={minimum height=10mm}]
\pic[val={B},] (b) at (4,-1.35) {2dmatrix};
\end{scope}
\begin{pgfonlayer}{background}\draw (du-out-0) to[out=0,in=180] (a-in-0);\end{pgfonlayer}
\begin{pgfonlayer}{background}\draw (du-out-1) to[out=0,in=180] (b-in-0);\end{pgfonlayer}
\begin{pgfonlayer}{background}\draw[bridge] (dv-out-0) to[out=0,in=180] (a-in-1);\end{pgfonlayer}
\begin{pgfonlayer}{background}\draw (dv-out-1) to[out=0,in=180] (b-in-1);\end{pgfonlayer}
\pic[] (gx) at (6.1,2.65) {ge};
\begin{pgfonlayer}{background}\draw (-.5,2.65) -- (gx-in-0);\end{pgfonlayer}
\begin{pgfonlayer}{background}\draw (gx-out-0) -- (7,2.65) to[out=0,in=0] (7,1.35) -- (a-out-0);\end{pgfonlayer}
\pic[] (gy) at (6.1,-1.35) {ge};
\begin{pgfonlayer}{background}\draw (b-out-0) -- (gy-in-0);\end{pgfonlayer}
\begin{pgfonlayer}{background}\draw (gy-out-0) -- (7,-1.35) to[out=0,in=0] (7,-2.65) -- (-.5,-2.65);\end{pgfonlayer}
\end{tikzpicture}}
\]


\[
\fits{\eqby{(C2)}\quad\begin{tikzpicture}[x=.77cm,y=.64cm,scale=0.88,baseline=(current bounding box.center)]
\pic[] (u) at (0,0.85) {codiscard};
\pic[yscale=.8] (du) at (1.3,0.85) {copy};
\begin{pgfonlayer}{background}\draw (u-out-0) -- (du-in-0);\end{pgfonlayer}
\pic[] (v) at (0,-0.85) {codiscard};
\pic[yscale=.8] (dv) at (1.3,-0.85) {copy};
\begin{pgfonlayer}{background}\draw (v-out-0) -- (dv-in-0);\end{pgfonlayer}
\node[font=\scriptsize] at (0,0) {$\vdots$};
\begin{scope}[matrixnode_left/.append style={minimum height=10mm}]
\pic[val={A},] (a) at (4,1.35) {2dmatrix};
\end{scope}
\begin{scope}[matrixnode_left/.append style={minimum height=10mm}]
\pic[val={B},] (b) at (4,-1.35) {2dmatrix};
\end{scope}
\begin{pgfonlayer}{background}\draw (du-out-0) to[out=0,in=180] (a-in-0);\end{pgfonlayer}
\begin{pgfonlayer}{background}\draw (du-out-1) to[out=0,in=180] (b-in-0);\end{pgfonlayer}
\begin{pgfonlayer}{background}\draw[bridge] (dv-out-0) to[out=0,in=180] (a-in-1);\end{pgfonlayer}
\begin{pgfonlayer}{background}\draw (dv-out-1) to[out=0,in=180] (b-in-1);\end{pgfonlayer}
\pic[] (sa) at (5.7,1.35) {antipode};
\pic[yscale=1.3] (sx) at (7.3,2) {multiplication};
\begin{pgfonlayer}{background}\draw (-.5,2.65) -- (sx-in-0);\end{pgfonlayer}
\begin{pgfonlayer}{background}\draw (a-out-0) -- (sa-in-0);\end{pgfonlayer}
\begin{pgfonlayer}{background}\draw (sa-out-0) -- (sx-in-1);\end{pgfonlayer}
\pic[] (sb) at (5.7,-2.65) {antipode};
\pic[yscale=1.3] (sy) at (7.3,-2) {multiplication};
\begin{pgfonlayer}{background}\draw (b-out-0) -- (sy-in-0);\end{pgfonlayer}
\begin{pgfonlayer}{background}\draw (-.5,-2.65) -- (sb-in-0);\end{pgfonlayer}
\begin{pgfonlayer}{background}\draw (sb-out-0) -- (sy-in-1);\end{pgfonlayer}
\pic[] (gx) at (8.7,2) {ge};
\pic[] (zx) at (10,2) {cozero};
\begin{pgfonlayer}{background}\draw (sx-out-0) -- (gx-in-0);\end{pgfonlayer}
\begin{pgfonlayer}{background}\draw (gx-out-0) -- (zx-in-0);\end{pgfonlayer}
\pic[] (gy) at (8.7,-2) {ge};
\pic[] (zy) at (10,-2) {cozero};
\begin{pgfonlayer}{background}\draw (sy-out-0) -- (gy-in-0);\end{pgfonlayer}
\begin{pgfonlayer}{background}\draw (gy-out-0) -- (zy-in-0);\end{pgfonlayer}
\end{tikzpicture}}
\]




\[
\fits{\eqby{(matrix def.), (row), (comp.)}\quad\begin{tikzpicture}[x=.77cm,y=.64cm,scale=0.9,baseline=(current bounding box.center)]
\begin{scope}[matrixnode_left/.append style={minimum height=24mm}]
\pic[val={M},] (m) at (3.2,0) {matrix};
\end{scope}
\pic[] (t) at (-0.5,1.55) {codiscard};
\begin{pgfonlayer}{background}\draw (t-out-0) -- ($(m-node.center)+(0,1.55)$);\end{pgfonlayer}
\pic[] (u) at (-0.5,0.7) {codiscard};
\begin{pgfonlayer}{background}\draw (u-out-0) -- ($(m-node.center)+(0,0.7)$);\end{pgfonlayer}
\pic[] (v) at (-0.5,0) {codiscard};
\begin{pgfonlayer}{background}\draw (v-out-0) -- ($(m-node.center)+(0,0)$);\end{pgfonlayer}
\node[font=\scriptsize] at (-0.85,0.35) {$\vdots$};
\begin{pgfonlayer}{background}\draw (-1,-.8) -- ($(m-node.center)+(0,-.8)$);\end{pgfonlayer}
\begin{pgfonlayer}{background}\draw (-1,-1.45) -- ($(m-node.center)+(0,-1.45)$);\end{pgfonlayer}
\pic[] (g) at (6.05,0) {ge};
\pic[] (z) at (7.35,0) {cozero};
\begin{pgfonlayer}{background}\draw (m-out-0) -- (g-in-0);\end{pgfonlayer}
\begin{pgfonlayer}{background}\draw (g-out-0) -- (z-in-0);\end{pgfonlayer}
\end{tikzpicture}}
\]

Separate one source.


\[
\fits{\eqby{(row)}\quad\begin{tikzpicture}[x=.77cm,y=.64cm,scale=0.88,baseline=(current bounding box.center)]
\begin{scope}[matrixnode_left/.append style={minimum height=17mm}]
\pic[val={R},] (r) at (2.7,0) {matrix};
\end{scope}
\pic[] (u) at (-0.5,0.7) {codiscard};
\pic[] (v) at (-0.5,0) {codiscard};
\begin{pgfonlayer}{background}\draw (u-out-0) -- ($(r-node.center)+(0,.7)$);\end{pgfonlayer}
\begin{pgfonlayer}{background}\draw (v-out-0) -- (r-in-0);\end{pgfonlayer}
\node[font=\scriptsize] at (-0.85,0.35) {$\vdots$};
\begin{pgfonlayer}{background}\draw (-1,-.65) -- ($(r-node.center)+(0,-.65)$);\end{pgfonlayer}
\begin{pgfonlayer}{background}\draw (-1,-1.25) -- ($(r-node.center)+(0,-1.25)$);\end{pgfonlayer}
\pic[] (t) at (-0.5,-2.8) {codiscard};
\pic[val={c},] (c) at (2.7,-2.8) {matrix};
\begin{pgfonlayer}{background}\draw (t-out-0) -- (c-in-0);\end{pgfonlayer}
\pic[yscale=2.8] (s) at (5.8,-1.4) {multiplication};
\begin{pgfonlayer}{background}\draw (r-out-0) -- (s-in-0);\end{pgfonlayer}
\begin{pgfonlayer}{background}\draw (c-out-0) -- (s-in-1);\end{pgfonlayer}
\pic[] (g) at (7.2,-1.4) {ge};
\pic[] (z) at (8.5,-1.4) {cozero};
\begin{pgfonlayer}{background}\draw (s-out-0) -- (g-in-0);\end{pgfonlayer}
\begin{pgfonlayer}{background}\draw (g-out-0) -- (z-in-0);\end{pgfonlayer}
\end{tikzpicture}}
\]


\[
\fits{\eqby{(C2)}\quad\begin{tikzpicture}[x=.77cm,y=.64cm,scale=0.88,baseline=(current bounding box.center)]
\begin{scope}[matrixnode_left/.append style={minimum height=17mm}]
\pic[val={R},] (r) at (2.7,0) {matrix};
\end{scope}
\pic[] (u) at (-0.5,0.7) {codiscard};
\pic[] (v) at (-0.5,0) {codiscard};
\begin{pgfonlayer}{background}\draw (u-out-0) -- ($(r-node.center)+(0,.7)$);\end{pgfonlayer}
\begin{pgfonlayer}{background}\draw (v-out-0) -- (r-in-0);\end{pgfonlayer}
\node[font=\scriptsize] at (-0.85,0.35) {$\vdots$};
\begin{pgfonlayer}{background}\draw (-1,-.65) -- ($(r-node.center)+(0,-.65)$);\end{pgfonlayer}
\begin{pgfonlayer}{background}\draw (-1,-1.25) -- ($(r-node.center)+(0,-1.25)$);\end{pgfonlayer}
\pic[] (t) at (-0.5,-2.8) {codiscard};
\pic[val={c},] (c) at (2.7,-2.8) {matrix};
\begin{pgfonlayer}{background}\draw (t-out-0) -- (c-in-0);\end{pgfonlayer}
\pic[] (a) at (4.35,-2.8) {antipode};
\begin{pgfonlayer}{background}\draw (c-out-0) -- (a-in-0);\end{pgfonlayer}
\pic[] (g) at (5,0) {ge};
\begin{pgfonlayer}{background}\draw (r-out-0) -- (g-in-0);\end{pgfonlayer}
\begin{pgfonlayer}{background}\draw (g-out-0) -- (6.5,0) to[out=0,in=0,looseness=1.3] (6.5,-2.8) -- (a-out-0);\end{pgfonlayer}
\end{tikzpicture}}
\]


\[
\fits{\eqby{(comp.)}\quad\begin{tikzpicture}[x=.77cm,y=.64cm,scale=0.88,baseline=(current bounding box.center)]
\begin{scope}[matrixnode_left/.append style={minimum height=17mm}]
\pic[val={R},] (r) at (2.7,0) {matrix};
\end{scope}
\pic[] (u) at (-0.5,0.7) {codiscard};
\pic[] (v) at (-0.5,0) {codiscard};
\begin{pgfonlayer}{background}\draw (u-out-0) -- ($(r-node.center)+(0,.7)$);\end{pgfonlayer}
\begin{pgfonlayer}{background}\draw (v-out-0) -- (r-in-0);\end{pgfonlayer}
\node[font=\scriptsize] at (-0.85,0.35) {$\vdots$};
\begin{pgfonlayer}{background}\draw (-1,-.65) -- ($(r-node.center)+(0,-.65)$);\end{pgfonlayer}
\begin{pgfonlayer}{background}\draw (-1,-1.25) -- ($(r-node.center)+(0,-1.25)$);\end{pgfonlayer}
\pic[] (t) at (-0.5,-2.8) {codiscard};
\pic[val={F},] (c) at (2.7,-2.8) {matrix};
\begin{pgfonlayer}{background}\draw (t-out-0) -- (c-in-0);\end{pgfonlayer}
\pic[] (g) at (5,0) {ge};
\begin{pgfonlayer}{background}\draw (r-out-0) -- (g-in-0);\end{pgfonlayer}
\begin{pgfonlayer}{background}\draw (g-out-0) -- (6.5,0) to[out=0,in=0,looseness=1.3] (6.5,-2.8) -- (c-out-0);\end{pgfonlayer}
\end{tikzpicture}}
\]


\[
\fits{\eqby{(slide)}\quad\begin{tikzpicture}[x=.77cm,y=.64cm,scale=0.88,baseline=(current bounding box.center)]
\begin{scope}[matrixnode_left/.append style={minimum height=17mm}]
\pic[val={R},] (r) at (2.7,0) {matrix};
\end{scope}
\pic[] (u) at (-0.5,0.7) {codiscard};
\pic[] (v) at (-0.5,0) {codiscard};
\begin{pgfonlayer}{background}\draw (u-out-0) -- ($(r-node.center)+(0,.7)$);\end{pgfonlayer}
\begin{pgfonlayer}{background}\draw (v-out-0) -- (r-in-0);\end{pgfonlayer}
\node[font=\scriptsize] at (-0.85,0.35) {$\vdots$};
\begin{pgfonlayer}{background}\draw (-1,-.65) -- ($(r-node.center)+(0,-.65)$);\end{pgfonlayer}
\begin{pgfonlayer}{background}\draw (-1,-1.25) -- ($(r-node.center)+(0,-1.25)$);\end{pgfonlayer}
\pic[] (g) at (5,0) {ge};
\pic[val={F},] (f) at (7.05,0) {comatrix};
\pic[] (e) at (9.1,0) {discard};
\begin{pgfonlayer}{background}\draw (r-out-0) -- (g-in-0);\end{pgfonlayer}
\begin{pgfonlayer}{background}\draw (g-out-0) -- (f-in-0);\end{pgfonlayer}
\begin{pgfonlayer}{background}\draw (f-out-0) -- (e-in-0);\end{pgfonlayer}
\node[draw,dashed,rounded corners=2pt,fit=(g-in-0)(g-node)(f-node)(e-node),inner sep=5pt] {};
\end{tikzpicture}}
\]


\[
\fits{\eqby{one-wire case; (matrix def.)}\quad\begin{tikzpicture}[x=.77cm,y=.64cm,scale=0.88,baseline=(current bounding box.center)]
\begin{scope}[matrixnode_left/.append style={minimum height=17mm}]
\pic[val={R},] (r) at (2.7,0) {matrix};
\end{scope}
\pic[] (u) at (-0.5,0.7) {codiscard};
\pic[] (v) at (-0.5,0) {codiscard};
\begin{pgfonlayer}{background}\draw (u-out-0) -- ($(r-node.center)+(0,.7)$);\end{pgfonlayer}
\begin{pgfonlayer}{background}\draw (v-out-0) -- (r-in-0);\end{pgfonlayer}
\node[font=\scriptsize] at (-0.85,0.35) {$\vdots$};
\begin{pgfonlayer}{background}\draw (-1,-.65) -- ($(r-node.center)+(0,-.65)$);\end{pgfonlayer}
\begin{pgfonlayer}{background}\draw (-1,-1.25) -- ($(r-node.center)+(0,-1.25)$);\end{pgfonlayer}
\pic[val={E},] (e) at (5.3,0) {matrix};
\pic[] (g) at (7.3,0) {ge};
\pic[] (z) at (8.6,0) {cozero};
\begin{pgfonlayer}{background}\draw (r-out-0) -- (e-in-0);\end{pgfonlayer}
\begin{pgfonlayer}{background}\draw (e-out-0) -- (g-in-0);\end{pgfonlayer}
\begin{pgfonlayer}{background}\draw (g-out-0) -- (z-in-0);\end{pgfonlayer}
\end{tikzpicture}}
\]


\[
\fits{\eqby{(comp.)}\quad\begin{tikzpicture}[x=.77cm,y=.64cm,scale=0.9,baseline=(current bounding box.center)]
\begin{scope}[matrixnode_left/.append style={minimum height=24mm}]
\pic[val={M'},] (m) at (3.2,0) {matrix};
\end{scope}
\pic[] (u) at (-0.5,0.7) {codiscard};
\begin{pgfonlayer}{background}\draw (u-out-0) -- ($(m-node.center)+(0,0.7)$);\end{pgfonlayer}
\pic[] (v) at (-0.5,0) {codiscard};
\begin{pgfonlayer}{background}\draw (v-out-0) -- ($(m-node.center)+(0,0)$);\end{pgfonlayer}
\node[font=\scriptsize] at (-0.85,0.35) {$\vdots$};
\begin{pgfonlayer}{background}\draw (-1,-.8) -- ($(m-node.center)+(0,-.8)$);\end{pgfonlayer}
\begin{pgfonlayer}{background}\draw (-1,-1.45) -- ($(m-node.center)+(0,-1.45)$);\end{pgfonlayer}
\pic[] (g) at (6.05,0) {ge};
\pic[] (z) at (7.35,0) {cozero};
\begin{pgfonlayer}{background}\draw (m-out-0) -- (g-in-0);\end{pgfonlayer}
\begin{pgfonlayer}{background}\draw (g-out-0) -- (z-in-0);\end{pgfonlayer}
\end{tikzpicture}}
\]



Repeat the same local rewrite until no source remains.


\[
\fits{\begin{tikzpicture}[x=.77cm,y=.64cm,scale=0.9,baseline=(current bounding box.center)]
\begin{scope}[matrixnode_left/.append style={minimum height=24mm}]
\pic[val={M'},] (m) at (3.2,0) {matrix};
\end{scope}
\pic[] (u) at (-0.5,0.7) {codiscard};
\begin{pgfonlayer}{background}\draw (u-out-0) -- ($(m-node.center)+(0,0.7)$);\end{pgfonlayer}
\pic[] (v) at (-0.5,0) {codiscard};
\begin{pgfonlayer}{background}\draw (v-out-0) -- ($(m-node.center)+(0,0)$);\end{pgfonlayer}
\node[font=\scriptsize] at (-0.85,0.35) {$\vdots$};
\begin{pgfonlayer}{background}\draw (-1,-.8) -- ($(m-node.center)+(0,-.8)$);\end{pgfonlayer}
\begin{pgfonlayer}{background}\draw (-1,-1.45) -- ($(m-node.center)+(0,-1.45)$);\end{pgfonlayer}
\pic[] (g) at (6.05,0) {ge};
\pic[] (z) at (7.35,0) {cozero};
\begin{pgfonlayer}{background}\draw (m-out-0) -- (g-in-0);\end{pgfonlayer}
\begin{pgfonlayer}{background}\draw (g-out-0) -- (z-in-0);\end{pgfonlayer}
\end{tikzpicture}\eqby{same rewrite}\ \cdots\ \eqby{same rewrite}\begin{tikzpicture}[x=.77cm,y=.64cm,scale=0.9,baseline=(current bounding box.center)]
\begin{scope}[matrixnode_left/.append style={minimum height=14mm}]
\pic[val={N},] (m) at (0,0) {matrix};
\end{scope}
\pic[] (t) at (-2.2,0.7) {codiscard};
\begin{pgfonlayer}{background}\draw (t-out-0) -- ($(m-node.center)+(0,.7)$);\end{pgfonlayer}
\begin{pgfonlayer}{background}\draw (-2.2,0) -- (m-in-0);\end{pgfonlayer}
\begin{pgfonlayer}{background}\draw (-2.2,-.7) -- ($(m-node.center)+(0,-.7)$);\end{pgfonlayer}
\pic[] (g) at (2.4,0) {ge};
\pic[] (z) at (3.7,0) {cozero};
\begin{pgfonlayer}{background}\draw (m-out-0) -- (g-in-0);\end{pgfonlayer}
\begin{pgfonlayer}{background}\draw (g-out-0) -- (z-in-0);\end{pgfonlayer}
\end{tikzpicture}}
\]


\[
\fits{\eqby{same rewrite}\quad\begin{tikzpicture}[x=.77cm,y=.64cm,scale=1,baseline=(current bounding box.center)]
\pic[val={C},] (m) at (0,0) {2dmatrix};
\pic[] (g) at (2,0) {ge};
\pic[] (z) at (3.3,0) {cozero};
\begin{pgfonlayer}{background}\draw (m-out-0) -- (g-in-0);\end{pgfonlayer}
\begin{pgfonlayer}{background}\draw (g-out-0) -- (z-in-0);\end{pgfonlayer}
\end{tikzpicture}}
\]

Apply the boundary-restoration chain above:


\[
\fits{\begin{tikzpicture}[x=.77cm,y=.64cm,scale=1,baseline=(current bounding box.center)]
\pic[val={C},] (m) at (0,0) {2dmatrix};
\pic[] (g) at (2,0) {ge};
\pic[] (z) at (3.3,0) {cozero};
\begin{pgfonlayer}{background}\draw (m-out-0) -- (g-in-0);\end{pgfonlayer}
\begin{pgfonlayer}{background}\draw (g-out-0) -- (z-in-0);\end{pgfonlayer}
\end{tikzpicture}
\eqby{cone form}
\begin{tikzpicture}[x=.77cm,y=.64cm,scale=1,baseline=(current bounding box.center)]
\pic[val={A'},] (a) at (0,0.75) {matrix};
\pic[val={B'},] (b) at (0,-0.75) {matrix};
\pic[] (g) at (2,0.75) {ge};
\begin{pgfonlayer}{background}\draw (a-out-0) -- (g-in-0);\end{pgfonlayer}
\begin{pgfonlayer}{background}\draw (g-out-0) -- (3.1,.75) to[out=0,in=0,looseness=1.3] (3.1,-.75) -- (b-out-0);\end{pgfonlayer}
\end{tikzpicture}}
\]


\[
\fits{\begin{tikzpicture}[x=.77cm,y=.64cm,scale=1,baseline=(current bounding box.center)]
        \pic[] (g) at (0,0) {ge};
        \pic[val={A},] (a) at (1.95,0) {comatrix};
        \pic[val={B},] (b) at (4.55,0) {matrix};
        \pic[] (h) at (6.5,0) {ge};
        \draw (g-out-0) -- (a-in-0);
        \draw (a-out-0) -- (b-in-0);
        \draw (b-out-0) -- (h-in-0);

        \end{tikzpicture}
\eqby{unbend}
\begin{tikzpicture}[x=.77cm,y=.64cm,scale=1,baseline=(current bounding box.center)]
        \pic[val={A'},] (a) at (0,0) {matrix};
        \pic[] (g) at (2,0) {ge};
        \pic[val={B'},] (b) at (4,0) {comatrix};
        \draw (a-out-0) -- (g-in-0);
        \draw (g-out-0) -- (b-in-0);

    \end{tikzpicture}}
\]

\end{proof}

\section{First normal form}\label{appendix:gpa-proof-5}

\begin{proof}
    We proceed by structural induction. The two base cases are as follows.
    First, the $\ge$ generator:
    \begin{center}
    \begin{tikzpicture}[axpic]
        \pic (g1) at (0.500,0.500) {ge};
    \end{tikzpicture}
    \;\begin{tikzpicture}[axpic]\node{$=$};\end{tikzpicture}\;
    \begin{tikzpicture}[axpic]
        \pic (g1) at (1.750,0.500) {ge};
        \coordinate (id_in_1)  at (0.250,0.500);
        \coordinate (id_out_1) at (1.250,0.500);
        \draw (id_in_1) -- (id_out_1);
        \coordinate (id_in_2)  at (2.250,0.500);
        \coordinate (id_out_2) at (3.250,0.500);
        \draw (id_in_2) -- (id_out_2);
        \draw[dashed] (0.350,0.150) rectangle (1.150,0.850);
        \draw[dashed] (2.350,0.150) rectangle (3.150,0.850);
    \end{tikzpicture}
    \;\begin{tikzpicture}[axpic]\node{$=$};\end{tikzpicture}\;
    \begin{tikzpicture}[axpic]
        \pic[val=I] (g1) at (0.500,0.500) {matrix};
        \begin{scope}[shift={(1.250,0.000)}]
            \pic (g2) at (0.500,0.500) {ge};
        \end{scope}
        \draw (g1-out-0) to[out=0,in=180] (g2-in-0);
        \begin{scope}[shift={(2.500,0.000)}]
            \pic[val=I] (g3) at (0.500,0.500) {comatrix};
        \end{scope}
        \draw (g2-out-0) to[out=0,in=180] (g3-in-0);
    \end{tikzpicture}
\end{center}


    \noindent Second, for any diagram $c$ in $\mathbf{GLA}$:
    \begin{center}
    \begin{tikzpicture}[axpic]
        \pic[val=A] (g1) at (0.500,0.500) {matrix};
    \end{tikzpicture}
    \;\begin{tikzpicture}[axpic]\node{$=$};\end{tikzpicture}\;
    \begin{tikzpicture}[axpic]
        \pic[val=A] (g1) at (0.500,1.000) {matrix};
        \pic (g2) at (1.750,1.000) {copy};
        \pic (g3) at (3.000,1.500) {ge};
        \pic (g4) at (3.000,0.500) {coge};
        \pic (g5) at (4.250,1.000) {cocopy};
        \draw (g1-out-0) to[out=0,in=180] (g2-in-0);
        \draw (g2-out-0) to[out=0,in=180] (g3-in-0);
        \draw (g2-out-1) to[out=0,in=180] (g4-in-0);
        \draw (g3-out-0) to[out=0,in=180] (g5-in-0);
        \draw (g4-out-0) to[out=0,in=180] (g5-in-1);
    \end{tikzpicture} \\
    \vspace{3mm}

    \;\begin{tikzpicture}[axpic]\node{$=$};\end{tikzpicture}\;
    \begin{tikzpicture}[axpic]
        \pic (g1) at (0.500,1.000) {copy};
        \pic[val=A] (g2) at (1.750,1.500) {matrix};
        \pic (g3) at (3.000,1.500) {ge};
        \pic[val=A] (g4) at (1.750,0.500) {matrix};
        \pic (g5) at (3.000,0.500) {coge};
        \pic (g6) at (4.250,1.000) {cocopy};
        \draw (g1-out-0) to[out=0,in=180] (g2-in-0);
        \draw (g1-out-1) to[out=0,in=180] (g4-in-0);
        \draw (g2-out-0) to[out=0,in=180] (g3-in-0);
        \draw (g4-out-0) to[out=0,in=180] (g5-in-0);
        \draw (g3-out-0) to[out=0,in=180] (g6-in-0);
        \draw (g5-out-0) to[out=0,in=180] (g6-in-1);
    \end{tikzpicture}
    \;\begin{tikzpicture}[axpic]\node{$\overset{(\text{antipode})}{=}$};\end{tikzpicture}\;
    \begin{tikzpicture}[axpic]
        \pic (g1) at (0.500,1.000) {copy};
        \pic[val=A] (g2) at (1.750,1.500) {matrix};
        \pic (g3) at (3.000,1.500) {ge};
        \pic[val=-A] (g4) at (1.750,0.500) {matrix};
        \pic (g5) at (3.000,0.500) {ge};
        \pic (g6) at (4.250,1.000) {cocopy};
        \draw (g1-out-0) to[out=0,in=180] (g2-in-0);
        \draw (g1-out-1) to[out=0,in=180] (g4-in-0);
        \draw (g2-out-0) to[out=0,in=180] (g3-in-0);
        \draw (g4-out-0) to[out=0,in=180] (g5-in-0);
        \draw (g3-out-0) to[out=0,in=180] (g6-in-0);
        \draw (g5-out-0) to[out=0,in=180] (g6-in-1);
    \end{tikzpicture} \\
    \vspace{3mm}

    \;\begin{tikzpicture}[axpic]\node{$\overset{(\text{matrix def.})}{=}$};\end{tikzpicture}\;
    \begin{tikzpicture}[axpic]
        \pic[val=\bar{A}] (g1) at (0.500,0.500) {matrix};
        \begin{scope}[shift={(1.250,0.000)}]
            \pic (g2) at (0.500,0.500) {ge};
        \end{scope}
        \draw (g1-out-0) to[out=0,in=180] (g2-in-0);
        \begin{scope}[shift={(2.500,0.000)}]
            \pic[val=\bar{B}] (g3) at (0.500,0.500) {comatrix};
        \end{scope}
        \draw (g2-out-0) to[out=0,in=180] (g3-in-0);
    \end{tikzpicture}
\end{center}


    \noindent
    For the inductive step, we consider composition of two diagrams
    $c; d$ in $\mathbf{GPA}_{\rtop}$. By the inductive hypothesis, both
    $c$ and $d$ are already in normal form:
    \begin{center}
    \begin{tikzpicture}[axpic]
        \pic[val=A] (g1) at (0.500,0.500) {matrix};
        \pic (g2) at (1.750,0.500) {ge};
        \pic[val=B] (g3) at (3.000,0.500) {comatrix};
        \draw (g1-out-0) to[out=0,in=180] (g2-in-0);
        \draw (g2-out-0) to[out=0,in=180] (g3-in-0);
    \end{tikzpicture}
    \;\begin{tikzpicture}[axpic]\node{$;$};\end{tikzpicture}\;
    \begin{tikzpicture}[axpic]
        \pic[val=C] (h1) at (0.500,0.500) {matrix};
        \pic (h2) at (1.750,0.500) {ge};
        \pic[val=D] (h3) at (3.000,0.500) {comatrix};
        \draw (h1-out-0) to[out=0,in=180] (h2-in-0);
        \draw (h2-out-0) to[out=0,in=180] (h3-in-0);
    \end{tikzpicture} \\
    \vspace{3mm}

    \;\begin{tikzpicture}[axpic]\node{$=$};\end{tikzpicture}\;
    \begin{tikzpicture}[axpic]
        \pic[val=A] (g1) at (0.500,0.500) {matrix};
        \pic (g2) at (1.750,0.500) {ge};
        \pic[val=B] (g3) at (3.000,0.500) {comatrix};
        \pic[val=C] (g4) at (4.250,0.500) {matrix};
        \pic (g5) at (5.500,0.500) {ge};
        \pic[val=D] (g6) at (6.750,0.500) {comatrix};
        \draw (g1-out-0) to[out=0,in=180] (g2-in-0);
        \draw (g2-out-0) to[out=0,in=180] (g3-in-0);
        \draw (g3-out-0) to[out=0,in=180] (g4-in-0);
        \draw (g4-out-0) to[out=0,in=180] (g5-in-0);
        \draw (g5-out-0) to[out=0,in=180] (g6-in-0);
        \draw[dashed] (1.150,0.050) rectangle (6.100,0.950);
    \end{tikzpicture} \\
    \vspace{3mm}

    \;\begin{tikzpicture}[axpic]\node{$\overset{(FME)}{=}$};\end{tikzpicture}\;
    \begin{tikzpicture}[axpic]
        \pic[val=\hat{A}] (g1) at (0.500,0.500) {matrix};
        \begin{scope}[shift={(1.250,0.000)}]
            \pic (g2) at (0.500,0.500) {ge};
        \end{scope}
        \draw (g1-out-0) to[out=0,in=180] (g2-in-0);
        \begin{scope}[shift={(2.500,0.000)}]
            \pic[val=\hat{D}] (g3) at (0.500,0.500) {comatrix};
        \end{scope}
        \draw (g2-out-0) to[out=0,in=180] (g3-in-0);
    \end{tikzpicture}
\end{center}


    \noindent For the tensor-product case:
    \begin{center}
    \begin{tikzpicture}[axpic]
        \pic[val=A] (g1) at (0.500,1.500) {matrix};
        \pic (g2) at (1.750,1.500) {ge};
        \pic[val=B] (g3) at (3.000,1.500) {comatrix};
        \draw (g1-out-0) to[out=0,in=180] (g2-in-0);
        \draw (g2-out-0) to[out=0,in=180] (g3-in-0);

        \pic[val=C] (h1) at (0.500,0.500) {matrix};
        \pic (h2) at (1.750,0.500) {ge};
        \pic[val=D] (h3) at (3.000,0.500) {comatrix};
        \draw (h1-out-0) to[out=0,in=180] (h2-in-0);
        \draw (h2-out-0) to[out=0,in=180] (h3-in-0);
    \end{tikzpicture}
    \;\begin{tikzpicture}[axpic]\node{$=$};\end{tikzpicture}\;
    \begin{tikzpicture}[axpic]
        \pic[val=\bar{A}] (g1) at (0.500,0.500) {matrix};
        \begin{scope}[shift={(1.250,0.000)}]
            \pic (g2) at (0.500,0.500) {ge};
        \end{scope}
        \draw (g1-out-0) to[out=0,in=180] (g2-in-0);
        \begin{scope}[shift={(2.500,0.000)}]
            \pic[val=\bar{B}] (g3) at (0.500,0.500) {comatrix};
        \end{scope}
        \draw (g2-out-0) to[out=0,in=180] (g3-in-0);
    \end{tikzpicture}
\end{center}


    \noindent where $\overline{A} = [A \; C]$ and $\overline{B} = [B \; D]$
\end{proof}

\section{Second normal form}\label{appendix:gpa-proof-6}

\begin{proof}
    Apply the first normal form to $c^\circ$. Applying polarity again
    recovers $c$, since polarity is involutive and respects the
    equational theory. The polar of an inequality normal form is a
    finitely generated normal form, which gives the required result.
    See \cite[Proposition~17 and Theorem~18]{bonchi2021polyhedral}.
\end{proof}

\section{Affine normal form}\label{appendix:gpa-proof-11}

\begin{proof}
    The derivation proceeds directly. Let $c$ be a
    diagram of \textbf{GPA}. Then we can collect all $\rtop$ on the left such
    that $c'$ is a diagram of $\textbf{GPA}_{\rtop}$

    \begin{figure}[H]
        \centering
        \begin{tikzpicture}
            \begin{scope}[shift={(0.000,0.250)}]
                \pic[val=C] (g1) at (0.500,0.500) {diagram};
            \end{scope}
            \node (rel6) at (1.500,0.750) {$=$};
            \begin{scope}[shift={(1.500,0.000)}]
                \begin{scope}[shift={(0.500,1.000)}]
                    \coordinate (id_in_2) at (0.25,0.050);
                    \coordinate (id_out_3) at (0.500,0.050);
                    \draw (id_in_2) -- (id_out_3);
                \end{scope}
                \pic (g4) at (0.700,0.450) {affine};
                \begin{scope}[shift={(1.750,0.250)}]
                    \pic[val=C'] (g5) at (0.500,0.500) {2ddiagram};
                \end{scope}
                \draw (id_out_3) to[out=0,in=180] (g5-in-0);
                \draw (g4-out-0) to[out=0,in=180] (g5-in-1);
            \end{scope}
        \end{tikzpicture}
    \end{figure}

    From Theorem~\ref{thm:FirstNormalForm} we know that $c'$ admits a normal form in
    terms of matrices $A$ and $B$. Decomposing $A$ into columns gives:

    \begin{figure}[H]
        \centering
        \begin{tikzpicture}
            \begin{scope}[shift={(0.000,0.550)}]
                \coordinate (id_in_1) at (0,0.500);
                \coordinate (id_out_2) at (1.000,0.500);
                \draw (id_in_1) -- (id_out_2);
            \end{scope}
            \pic (g3) at (0.500,0.450) {affine};
            \begin{scope}[shift={(1.750,0.250)}]
                \pic[val=A] (g4) at (0.500,0.500) {2dmatrix};
            \end{scope}
            \draw (id_out_2) to[out=0,in=180] (g4-in-0);
            \draw (g3-out-0) to[out=0,in=180] (g4-in-1);
            \begin{scope}[shift={(3.500,0.250)}]
                \pic (g5) at (0.500,0.500) {ge};
            \end{scope}
            \draw (g4-out-0) to[out=0,in=180] (g5-in-0);
            \begin{scope}[shift={(5.250,0.250)}]
                \pic[val=B] (g6) at (0.500,0.500) {comatrix};
            \end{scope}
            \draw (g5-out-0) to[out=0,in=180] (g6-in-0);
            \node at (7,0.7) {$=$};
        \end{tikzpicture}\\
        \vspace{5mm}
        \begin{tikzpicture}
            \node at (-.5,1) {$=$};
            \begin{scope}[shift={(0.750,1.000)}]
                \pic[val=A_1] (g1) at (0.500,0.500) {matrix};
            \end{scope}
            \pic (g2) at (0.500,0.500) {affine};
            \begin{scope}[shift={(1.500,0.000)}]
                \pic[val=A_2] (g3) at (0.500,0.500) {matrix};
            \end{scope}
            \draw (g2-out-0) to[out=0,in=180] (g3-in-0);
            \coordinate (a4) at (0.000,0 |- g1-in-0);
            \draw (g1-in-0) -- (a4);
            \coordinate (a5) at (2.500,0 |- g1-out-0);
            \draw (g1-out-0) -- (a5);
            \begin{scope}[shift={(3.500,0.500)}]
                \pic (g6) at (0.500,0.500) {multiplication};
            \end{scope}
            \draw (a5) to[out=0,in=180] (g6-in-0);
            \draw (g3-out-0) to[out=0,in=180] (g6-in-1);
            \begin{scope}[shift={(5.500,0.500)}]
                \pic (g7) at (0.500,0.500) {ge};
            \end{scope}
            \draw (g6-out-0) to[out=0,in=180] (g7-in-0);
            \begin{scope}[shift={(7.500,0.500)}]
                \pic[val=B] (g8) at (0.500,0.500) {comatrix};
            \end{scope}
            \draw (g7-out-0) to[out=0,in=180] (g8-in-0);
        \end{tikzpicture}
    \end{figure}

    This proves the claimed normal form.
\end{proof}

